\chapter{Feldermodell}\label{chapter:Feldermodell}

In diesem Kapitel geht es um die lineare Stellung der unmittelbaren Satzbausteine in Bezug auf die Stellung der Verbteile. Der folgende Satz in (\ref{ea:Klammer_dt}) wird erst im Vergleich zur englischen Übersetzung (\ref{ex:KeineKlammer_engl}) auffällig. 

\ea \label{ea:Satzklammer}
\ea \label{ea:Klammer_dt} Anna wird morgen dem Jungen einen Ball schenken.
\ex \label{ex:KeineKlammer_engl} Tomorrow Anna will give the boy a ball.
\z 
\z

Im deutschen\is{Aussagesatz} Aussagesatz (\ref{ea:Klammer_dt}) stehen das finite Verb \textit{wird} und das infinite Verb \textit{schenken} voneinander getrennt: \textit{wird \ldots{} schenken}. Im Englischen (\ref{ex:KeineKlammer_engl})  -- wie auch in den romanischen Sprachen -- ist das nicht der Fall. Die zusammengehörigen verbalen Teile \textit{will give} stehen zusammen. Das Deutsche zeichnet sich durch die getrennte Stellung der Verbteile aus, die man als Klammer bezeichnet. Während im Englischen sowohl \textit{tomorrow} als auch \textit{Anna} vor dem finiten Verb \textit{will} stehen, ist das im Deutschen nicht möglich: *\textit{Morgen Anna will dem Jungen einen Ball schenken}.   

Das finite Verb muss jedoch nicht wie in (\ref{ea:Klammer_dt}) an der zweiten Stelle stehen. Es kann auch wie in (\ref{ea:ErstFin}) an der ersten Stelle oder wie in (\ref{exEndFIN}) ganz am Ende stehen. Bei dieser Stellung bilden die Verbteile keine Klammer.

\ea 
\ea \label{ea:ErstFin} Wird Anna dem Jungen einen Ball schenken?
\ex \label{exEndFIN} (Alle wissen,) dass Anna dem Jungen einen Ball schenken wird.
\z 
\z 

Je nach Stellung des finiten Verbs unterscheidet man im Deutschen zwischen Verbzweit-, Verberst- und Verbletztsätzen. Diesen drei Satztypen lassen sich Satzarten nach ihrer kommunikativen Funktion zuordnen. Hierum geht es im folgenden Abschnitt~\ref{sec:Satztypen_Satzarten}. Danach beschäftigen wir uns mit den einzelnen Satztypen, mit den in ihnen enthaltenen Feldern und ihrer Besetzung. In Abschnitt~\ref{section:Verbzweitsätze} geht es um Verbzweitsätze, in Abschnitt~\ref{section:Verberstsatz} um Verberstsätze und in Abschnitt~\ref{section:Verbletztsatz} um Verbletztsätze. Danach werden wir in Abschnitt~\ref{section:differenziert_uniform_Feldermodell} diskutieren, inwieweit wir für die verschiedenen Satztypen jeweils ein eigenes Modell benötigen oder alle Satztypen mithilfe eines Modells erfassen können. Abschließend erklären wir in Abschnitt~\ref{section:DeutschV2VletztfreieWortstellung}, inwiefern Deutsch eine Verbzweitsprache mit Verbletztstellung und freier Wortfolge ist. Stichpunktartig fassen wir dieses Kapitel im Abschnitt~\ref{section:Zfsg_Feldermodell} zusammen. Der letzte Abschnitt~\ref{section:Aufgaben_Feldermodell} bietet vier verschiedene Anwendungsaufgaben samt Lösungen zum Feldermodell. 


\section{Satztypen und Satzarten}\label{sec:Satztypen_Satzarten}

Bei komplexen Verben ist für einen deutschen\is{Aussagesatz} Aussagesatz wie (\ref{ea:Klammer_dt}) die Distanzstellung von finitem (\textit{wird}) und infinitem Verb (\textit{schenken}) wesentlich. Diese Distanzstellung bezeichnet man als\is{Verbklammer} \textit{Verbklammer}. Von der Verbklammer aus können\is{Feldermodell} wir bestimmte Felder bestimmen, auf denen die unmittelbaren Satzbausteine, die Konstituenten des Satzes stehen. Es handelt sich um Phrasen, die in erster Linie die Ergänzungen und Angaben des (Haupt-)Verbs realisieren. Vor dem Finitum gibt es das sogenannte Vorfeld und zwischen Finitum und Infinitiv das Mittelfeld. Für den Satz (\ref{ea:Klammer_dt}) kann das vorläufig wie in der Abbildung~\ref{fig:Felder_Aussagesatz} dargestellt werden.

\begin{figure}[!htbp]
  \centering
  \begin{tabular}{cp{0.1em}cp{0.1em}cp{0.1em}c}
    Vorfeld && Finitum && Mittelfeld && Infinitiv \\
    \cmidrule{1-1}\cmidrule{3-3}\cmidrule{5-5}\cmidrule{7-7}
    Anna && wird && dem Jungen einen Ball && schenken. \\
  \end{tabular}
  \caption{Felder im Aussagesatz (vorläufig)}
  \label{fig:Felder_Aussagesatz}
\end{figure}  

Wesentlich ist, dass das\is{Vorfeld} Vorfeld im deutschen Aussagesatz nur von einer Satzkonstituente besetzt werden kann. Deshalb ist die Stellung wie in (\ref{ea:doppVF}) ungrammatisch. Ungrammatisch ist auch die Stellung in (\ref{ex:Infnicht RSK}), wo die Verbteile keine Klammer bilden.

\ea \label{ea:VerstößeStellung}
\ea [*] {Anna dem Jungen wird einen Ball schenken.}\label{ea:doppVF}
\ex [*] {Anna wird schenken dem Jungen einen Ball.}\label{ex:Infnicht RSK}
\z 
\z 

In (\ref{ea:doppVF}) darf entweder nur \textit{Anna} oder nur \textit{dem Jungen} vor dem finiten Verb \textit{wird} stehen. In (\ref{ex:Infnicht RSK}) muss der Infinitiv \textit{schenken} am Ende stehen und somit das\is{Mittelfeld} Mittelfeld begrenzen. Im Mittelfeld können mehrere Satzkonstituenten wie \zB \textit{dem Jungen} und \textit{einen Ball} stehen.

Da die\is{Feldermodell} Feldstruktur durch die Stellung der Verbteile bestimmt wird, muss zwischen drei grundlegenden Stellungstypen als Formtypen unterschieden werden. Zentral für diese Unterscheidung ist die Stellung des finiten Verbs. Das Finitum kann wie in (\ref{ea:VerbzweitBsp}) an der zweiten Stelle stehen oder wie in (\ref{ex:VerberstBsp}) an der ersten Stelle oder wie in (\ref{ex:VerbendBsp}) an der letzten Stelle.   

\ea
\ea \label{ea:VerbzweitBsp} Anna \textit{wird} dem Jungen einen Ball schenken. (Verbzweit)
\ex \label{ex:VerberstBsp} \textit{Wird} Anna dem Jungen einen Ball schenken? (Verberst)
\ex \label{ex:VerbendBsp} dass Anna dem Jungen einen Ball schenken \textit{wird} (Verbletzt) 
\z 
\z 

Für diese drei formen Satztypen sind die Begriffe Verbzweit-, Verberst- und Verbletztsatz etabliert. Eigentlich müsste es Finitzweit-, Finiterst- und Finitletztsatz heißen.

Den formalen Satztypen kann man bestimmte funktionale Satzarten zuordnen. Der Prototyp\is{Verbzweitsatz} des Satzes ist der\is{Aussagesatz} Aussagesatz als eine Funktion des Verbzweitsatzes (\ref{ea:VerbzweitBsp}). Mit Aussagesätzen handeln wir, indem wir etwas als wahr behaupten. Erst mit einem Aussagesatz wird explizit eine Behauptung gemacht. Diese kann wahr oder falsch sein. Die zweite wichtige Funktion von Verbzweitsätzen besteht im Erfragen von Informationen mit W-Fragewörtern, \zB \textit{Was wird Anna dem Jungen schenken?}

Mit Verberstsätzen\is{Verberstsatz} (\ref{ex:VerberstBsp}) handeln wir auch, indem wir den Wahrheitswert der Behauptung erfragen. Verberstsätze werden auch imperativisch gebraucht wie \textit{Schenk ihm den Ball!} Dann fordern wir unser Gegenüber zu der vom Verb bedeuteten Handlung auf. Ihr Wahrheitswert hängt von ihrer Ausführung ab. 

Mit Verbletztsätzen\is{Verbletztsatz} (\ref{ex:VerbendBsp}) machen wir Behauptungen nur indirekt. Im Vergleich zum Behauptungsgehalt von Verberst- und Verbzweitsätzen sind Verbletztsätze neutral. Typischerweise fungieren Verbletztsätze als Ergänzung oder Angabe eines übergeordneten Satzes, der den Behauptungsgehalt bestimmt.   

Kommunikative Handlungen wie Behaupten, Fragen, Auffordern etc. nennt man\is{Illokution} \textit{Illokutionen}.\footnote{Die Illokution ist der wichtigste Sprechakt in der Sprechakttheorie von \citet{Austin1962} und \citet{Searle1969}. Man unterscheidet zwischen dem Hervorbringen einer sprachlichen Äußerung (lokutiver Akt), dem Zweck dieser Äußerung (illokutiver Akt) und der Wirkung, die mithilfe des illokutionären Aktes beim Adressaten bzw. der Adressatin erzielt wird (perlokutiver Akt). Einen kurzen Überblick bieten \citet[218--220]{Stukenbrock2013} und \citet[146--148]{Klabunde2023Pragmatik}.} Die Illokution liegt über der Prädikation (vgl. Definition~\ref{Def_Prädikation}).

\begin{Definition}{Illokution}
    
\noindent 
Illokution ist die sprachliche Handlung, die durch eine sprachliche Äußerung in einem kommunikativen Kontext vollzogen wird. Zentrale illokutionäre Handlungen sind das Behaupten, das Erfragen und das Auffordern. 
\label{Def_Illokution}
\end{Definition}



In den folgenden drei Abschnitten beschreiben wir die Verbzweit-, Verberst- und Verbletztsätze und beschäftigen uns mit der Besetzung ihrer Felder. Zuvor gehen wir in dem folgenden Exkurs~\ref{ExkursFeldermodellSchule} kurz auf den Feldbegriff und seine Rolle in der Schule ein. 

\begin{Exkurs}{Der Feldbegriff und die Schule}
\label{ExkursFeldermodellSchule}

Der \textit{Feld}begriff geht auf \citet[27]{Drach1963} (erstmals 1937 veröffentlicht) zurück. Die für die\is{Feldermodell} Felderbildung wesentlichen Nähe- und Distanzstellungen des finiten Verbs zum Verbrest hat \citet[298--306]{Herling1821} beschrieben. Durch \citet[310--313]{Becker1837} finden sie Eingang in die\is{Schulgrammatik} Schulgrammatik, sind aber bisher nur im Bereich Deutsch als Fremd- und Zweitsprache wirklich angekommen. \citet[181--197]{Erdmann1886} systematisiert die Satzfelder nach Satztypen in Abhängigkeit von der Stellung des Finitums an der zweiten, der ersten oder der letzten Stelle. Die in der Schule gängigen Begriffe \textit{Kernsatz} für Verbzweitsatz, \textit{Stirnsatz} für Verberstsatz und \textit{Spannsatz} für \textit{Verbletztsatz} gehen auf \citet[96]{Glinz1952} zurück.\footnote{Einen Forschungsrückblick und eine theoretische Diskussion der Felder bietet \citet{Hoehle19862018} in seiner 1986 erstmals veröffentlichten Darstellung. \citet[288]{Hoehle19862018} bezeichnet das Feldermodell auch als \textit{Herling/Erdmann-System}.} 

Während das Feldmodell im Fremdsprachenunterricht immer eine zentrale Rolle spielte, taucht es erst nach und nach wieder in Lehrbüchern und Materialien für den Deutschunterricht auf. Denn das Feldermodell ermöglicht es auf anschauliche Art und Weise, den linearen Aufbau von Sätzen zu reflektieren und in Abhängigkeit von der kommunikativen Absicht korrekt und variabel zu gestalten. Die Grundlage hierfür bilden die drei syntaktischen Grundmuster: der Verbzweit-, der Verberst- und Verbletztsatz mit entsprechenden Feldern und deren Besetzungsregeln. Den Nutzen des Feldermodells für den Unterricht fasst \citet[104]{Duerscheid2012} wie folgt zusammen: 

\begin{quote}
Die Untergliederung des Satzes in Stellungsfelder ermöglicht es, Wortstellungsregularitäten präziser zu formulieren und für jeden Verbstellungstyp die normalen von den abweichenden Felderbesetzungen zu erfassen. \citep[104]{Duerscheid2012} 
\end{quote}

Erst einmal unabhängig davon, wie das Modell didaktisch veranschaulicht wird: ob als Plätze in einem Bus,\footnote{So \zB in Metzger, Stefan (2021): \textit{Grammatik unterrichten mit dem Feldermodell. Didaktische Grundlagen und Aufgaben für die Orientierungsstufe}. 2. Auflage. Seelze: Klett, Kallmeyer.} mithilfe eines Klammermanns, dessen bewegliche Arme die Verbklammer bilden\footnote{So vorgeschlagen von \citet{Schoenenberg2011}.} oder ganz einfach mittels einer Tabelle an der Tafel, deren Spalten flexibel mit magnetischen Wort- oder Wortgruppenkarten besetzt werden -- zukünftige Lehrerinnen und Lehrer müssen verstehen, was dahintersteckt, und auch in der Lage sein, die entsprechenden didaktischen Materialien zu bewerten. Konkrete Vorschläge für den schulischen Einsatz des Feldermodells macht \citet[11]{Froemel2020} im Rahmen einer Lernprogression für die fünfte bis zur zwölften Klasse. \citet{Bredel2015TopOrth} stellt das Potenzial des Feldermodells für die Arbeit an der Orthografie und die Ermittlung ihres Erwerbs dar. Ihr geht es insbesondere um die Getrennt- und Zusammenschreibung sowie um die Zeichensetzung.

\end{Exkurs}



\section{Verbzweitsätze}\label{section:Verbzweitsätze}
\largerpage
Viele Verben sind im Deutschen zweiteilig. Sie bestehen aus einem einfachen Verb und einem weiteren Teil, \zB einem trennbaren Präfix (Verbpartikel) (\ref{ea:V2Pt}) oder einem Infinitiv (\ref{ex:V2Inf}) oder einem anderen in das Verb inkorporierten Teil (\ref{ex:V2Inkorp}). In Aussagesätzen steht das finite Verb\is{Verbzweitsatz} an zweiter Stelle und der Verbrest an letzter. Die Verbzweitstellung ist ein wesentliches Merkmal des Deutschen bzw. mit Ausnahme des Englischen der germanischen Sprachen.

\ea
\ea \label{ea:V2Pt} Anna \textit{ruft} mich morgen \textit{an}.
\ex \label{ex:V2Inf} Anna \textit{geht} gerne \textit{spazieren}.
\ex \label{ex:V2Inkorp} Anna \textit{steht} wieder \textit{kopf}. 
\z 
\z 

Auch lexikalisch einfache Verben wie \textit{schenken} werden \zB zum Ausdruck des Perfekts (\textit{hat\ldots{} geschenkt}) oder beim Gebrauch mit Modalverben (\textit{will\ldots{} schenken}) zweiteilig oder mehrteilig gebraucht (vgl. Abschnitt~\ref{section:RegierteVerben}). Auch in diesen Fällen steht die finite Form des Hilfs- oder Modalverbs an zweiter Stelle und der infinite Verbteil, das Partizip Perfekt (\ref{ea:hatgeschenkt}) oder der Infinitiv (\ref{ex:willschenken}) am Ende. Wenn es nur ein finites Verb ohne infiniten Verbteil gibt, so steht dieses an zweiter Stelle (\ref{ex:schenktV2}). Die jeweilige Besetzung der ersten Stelle ist hierbei unerheblich.

\ea 
\ea \label{ea:hatgeschenkt} Anna \textit{hat} dem Jungen einen Ball \textit{geschenkt}.
\ex \label{ex:willschenken} Einen Ball \textit{will} Anna dem Jungen \textit{schenken}.
\ex \label{ex:schenktV2} Dem Jungen \textit{schenkt} Anna einen Ball.
\z 
\z 

Das finite Verb und der infinite Verbteil bzw. der Verbrest bilden wie in den Beispielen (\ref{ea:V2Pt}) bis (\ref{ex:willschenken}) durch ihre Distanzstellung eine Klammer. Diese Klammer wird als\is{Satzklammer} \textit{Satzklammer} bezeichnet. Das Finitum bildet die\is{Satzklammer!linke Satzklammer} \textit{linke Satzklammer} (LSK).\footnote{Deshalb bezeichnet \citet[280]{Hoehle19862018} die LSK als \textit{FINIT}-Position.} Die infiniten Verbteile bzw. der Verbrest  bilden die\is{Satzklammer!rechte Satzklammer} \textit{rechte Satzklammer} (RSK). Zwischen linker und rechter Satzklammer, also innerhalb der Satzklammer, befindet sich das\is{Mittelfeld} \textit{Mittelfeld} (MF). Die Stelle vor der linken Satzklammer nennt man\is{Vorfeld} \textit{Vorfeld} (VF). Auch nach der rechten Klammer gibt es noch ein Feld, das sogenannte\is{Nachfeld} \textit{Nachfeld} (NF). Die folgende Abbildung~\ref{fig:FeldermodelV2voll} stellt die Feldbesetzung und Klammerbildung in den Sätzen (\ref{ea:V2Pt}) bis (\ref{ex:schenktV2}) dar. Zur Illustration für die mögliche Besetzung des Nachfelds haben wir einige Sätze ein bisschen ausgebaut. Zur Illustration einer komplexen rechten Satzklammer haben wir \textit{spazieren gehen} aus (\ref{ex:V2Inf}) ins Perfekt gesetzt. Die Striche kennzeichnen die Nicht-Besetzung vorhandener Felder bzw. der RSK. 

\begin{figure}[!htbp]
  \centering
  %\resizebox{\textwidth}{!}{
  \oneline{
    \begin{tabular}{cp{0.1em}cp{0.1em}cp{0.1em}cp{0.1em}c}
      VF && LSK && MF && RSK && NF \\
      \cmidrule{1-1}\cmidrule{3-3}\cmidrule{5-5}\cmidrule{7-7}\cmidrule{9-9}
      Anna && ruft && mich morgen && an, && wenn sie es nicht vergisst. \\
       Anna && ist && gestern && spazieren gegangen. && -- \\
        Anna && steht && wieder && kopf, && weil der Schlüssel weg ist. \\
      Anna && hat && dem Jungen einen Ball && geschenkt, && weil er Geburtstag hat.\\
      Einen Ball && will && Anna dem Jungen && schenken. && -- \\
       Dem Jungen && schenkt && Anna einen Ball, && -- && weil er Geburtstag hat. \\
      Anna && schläft. && -- && -- && -- \\ 
    \end{tabular}
  }
  \caption{Feldermodell für Verbzweitsätze}
  \label{fig:FeldermodelV2voll}
\end{figure}

In\is{Verbzweitsatz} Verbzweitsätzen sind nur VF und LSK obligatorisch zu besetzen. Die anderen Felder können valenzbedingt oder bei Einteiligkeit des Verbs unbesetzt bleiben, was anhand von \textit{Anna schläft} in der Abbildung~\ref{fig:FeldermodelV2voll} ersichtlich wird.

Aus fremd- und zweitsprachendidaktischer Perspektive regt \citet[75--76]{Dimroth2009} an, im Unterricht von Anfang an mit klammernden Prädikatsteilen zu arbeiten, damit die Kategorie der Finitheit in der LSK von der Kategorie des Vollverbs in der RSK getrennt wird. 

Verbzweitsätze bezeichnet man auch als \textit{Kernsätze}. Sie dienen primär dem Ausdruck von\is{Aussagesatz} Aussage- (\ref{ea:AussageV2}) und von\is{Ergänzungsfragesatz} Ergänzungsfragesätzen (W-Fragen) (\ref{ex:ErgFragenV2}). Mit ersteren treffen wir Behauptungen. Mit letzteren erfragen wir Informationen. Möglich sind ebenso Kernsätze zum Ausdruck einer Aufforderung (\ref{ex:AufforderungV2}) oder eines Ausrufs (\ref{ex:AusrufV2}).

\ea \label{ea:SatzartenV2}
\ea \label{ea:AussageV2} Anna schenkt dem Jungen einen Ball.
\ex \label{ex:ErgFragenV2} Was schenkt Anna dem Jungen?
\ex \label{ex:AufforderungV2} Du schenkst jetzt dem Jungen einen Ball!
\ex \label{ex:AusrufV2} Das war eine Freude!
\z 
\z 

Allerdings können Verbzweitsätze wie \textit{die Ministerin habe das Vertrauen des Kanzlers} auch als abhängige, also unselbstständige Nebensätze wie in (\ref{ea:V2_NeSa}) gebraucht werden. Diese uneingeleiteten Nebensätze stehen im Kontrast zu eingeleiteten Nebensätzen mit Verbletztstellung wie \textit{dass die Ministerin das Vertrauen des Kanzlers habe} in (\ref{ex:Vletzt_NeSa}). 

\ea \label{ea:V2vsVletzt}
\ea \label{ea:V2_NeSa} Ein Regierungssprecher betonte, die Ministerin habe das Vertrauen des Kanzlers.
\ex \label{ex:Vletzt_NeSa} Ein Regierungssprecher betonte, dass die Ministerin das Vertrauen des Kanzlers habe.
\z
\z 

Auf uneingeleitete Nebensätze mit Verbzweitstellung wie in (\ref{ea:V2_NeSa}) gehen wir in Abschnitt~\ref{subsection:FormFktNeSa} unter dem Stichwort \textit{verkappter Nebensatz} genauer ein. In den folgenden Unterkapiteln geben wir eine Übersicht über die Besetzungsmöglichkeiten der Felder von Verbzweitsätzen.

\subsection{Vorfeld}\label{subsection:Vorfeld}

Das VF\is{Vorfeld} muss von \textit{einem} unmittelbaren Satzbaustein besetzt sein. Um unmittelbare Satzbausteine geht es in Kapitel \ref{chap:Konstituenz} unter dem Begriff der unmittelbaren Satzkonstituenten. Typischerweise handelt es sich dabei um eine Wortgruppe, die \textit{eine} Ergänzung oder \textit{eine} Angabe des Verbs aus dem Hauptsatz realisiert.\footnote{Typischerweise, da auch Konstituenten vorangestellt werden können, die nicht direkt vom Hauptverb abhängen. Von Fällen wie \textit{Seiner Schwester soll der Präsident versucht haben, einen Ministerposten zu verschaffen} sehen wir hier ab und verweisen bei Interesse auf \citet[Kapitel 9]{Mueller1999}.} Dies gilt auch dann, wenn anstelle der jeweiligen Ergänzung ein Fragewort in sogenannten Ergänzungsfragen das VF besetzt, was die Abbildung~\ref{fig:FeldermodelV2mitW-Frage} zeigt.

\begin{figure}[!htbp]
  \centering
  %\resizebox{\textwidth}{!}{
  \oneline{
    \begin{tabular}{cp{0.1em}cp{0.1em}cp{0.1em}cp{0.1em}c}
      VF && LSK && MF && RSK && NF \\
      \cmidrule{1-1}\cmidrule{3-3}\cmidrule{5-5}\cmidrule{7-7}\cmidrule{9-9}
      Anna && will && dem Jungen einen Ball && schenken, && weil er Geburtstag hat. \\
      Wem && will && Anna einen Ball && schenken, && weil er Geburtstag hat? \\
       Was && will && Anna dem Jungen && schenken, && weil er Geburtstag hat? \\
       Warum && will && Anna dem Jungen einen Ball && schenken? && -- \\
    \end{tabular}
  }
  \caption{Besetzung des Vorfelds durch Fragepronomen}
  \label{fig:FeldermodelV2mitW-Frage}
\end{figure}

Auch Nebensätze können im\is{Vorfeld} VF stehen, wenn es sich um unmittelbare Satzkonstituenten handelt. So kann \zB die Angabe \textit{weil er Geburtstag hat} genauso im Vorfeld stehen wie Nebensätze und \textit{zu}-Infinitiv-Phrasen, die Ergänzungen des Verbs realisieren. Dies zeigen wir in der Abbildung~\ref{fig:FeldermodelV2mitNeSaVF}.

\begin{figure}[!htbp]
  \centering
  %\resizebox{\textwidth}{!}{
  \oneline{
    \begin{tabular}{cp{0.1em}cp{0.1em}cp{0.1em}cp{0.1em}c}
      VF && LSK && MF && RSK && NF \\
      \cmidrule{1-1}\cmidrule{3-3}\cmidrule{5-5}\cmidrule{7-7}\cmidrule{9-9}
      Weil er Geburtstag hat, && will && Anna dem Jungen einen Ball && schenken. && -- \\
      Dass du mir hilfst, && freut && mich. && -- && -- \\
       Dass du mir hilfst, && hast && du mir && versprochen. && -- \\
       Mir zu helfen && hast && du mir && versprochen. && -- \\ 
    \end{tabular}
  }
  \caption{Besetzung des Vorfelds durch Nebensätze und \textit{zu}-Infinitiv-Phrasen}
  \label{fig:FeldermodelV2mitNeSaVF}
\end{figure}

Zwei verschiedene Konstituenten, die nicht gemeinsam \textit{eine} Satzkonstituente bilden, dürfen das\is{Vorfeld} VF nicht gleichzeitig besetzen. Nicht vorfeldfähig sind auch Partikeln\footnote{Im Gegensatz zur Stellung im Mittelfeld: \textit{Du weißt das ja}.} (vgl. Abschnitt~\ref{subsection:Partikeln}) und Reflexivpronomen echt reflexiver Verben (vgl. Abschnitt~\ref{subsection:Reflexiva}). Dies wird in der Abbildung~\ref{fig:nichtVFfähig} gezeigt.\footnote{Eine Konstituente aus dem NF, die nicht im MF stehen kann, darf auch nicht das VF besetzen. \citet[291]{Hoehle19862018} zeigt dies in Bezug auf den Nebensatz in \textit{Wahrscheinlich ist Karl stolz gewesen, daß er so viel getrunken hatte} im Vergleich zu *\textit{Karl ist, daß er so viel getrunken hatte, stolz gewesen} und damit auch zu *\textit{Daß er so viel getrunken hatte, ist Karl stolz gewesen}.}   

\begin{figure}[!htbp]
  \centering
  %\resizebox{\textwidth}{!}{
  \oneline{
    \begin{tabular}{cp{0.1em}cp{0.1em}cp{0.1em}cp{0.1em}c}
      VF && LSK && MF && RSK && NF \\
      \cmidrule{1-1}\cmidrule{3-3}\cmidrule{5-5}\cmidrule{7-7}\cmidrule{9-9}
      * Anna weil er Geburtstag hat && will && dem Jungen einen Ball && schenken. && -- \\
       * Anna dem Jungen && will && einen Ball && schenken. && -- \\
       * Ja && weißt && du das. && -- && --\\
       * Sich && schämt && Paul. && -- && -- 
    \end{tabular}
  }
  \caption{Unmögliche Vorfeldbesetzung}
  \label{fig:nichtVFfähig}
\end{figure}

Bei mehrteiligen Verbformen kann auch der infinite Verbteil aus der RSK im\is{Vorfeld} VF stehen. Der Infinitiv oder das Partizip II kann dabei als Valenzträger seine Ergänzungen und/oder Angaben mit ins VF nehmen. Die erste verbale Ergänzung, in der Regel das Subjekt, ist normalerweise von der Voranstellung mit dem infiniten Prädikatsteil ausgeschlossen.\footnote{Nur bei intransitiven Verben, deren erste Ergänzung das Geschehen nicht kontrolliert und die ihr Perfekt mit \textit{sein} bilden, kann die Nominativ-Ergänzung zusammen mit dem Partizip II im VF stehen, insbesondere dann, wenn es noch eine belebte Dativ-NP gibt, \zB \textit{Ein Fehler unterlaufen ist ihm noch nie}. Zur Diskussion verweisen wir auf \citet{Fanselow1992}.} Die Stellung infiniter Verbteile im VF wird in der folgenden Abbildung~\ref{fig:FeldermodelV2mitINFimVF} dargestellt. Das NF bleibt der Übersichtlichkeit halber unbesetzt.\footnote{Ebenso sehen wir von einer möglichen Besetzung der RSK ab. Möglich wäre \zB \textit{Einen Ball schenken will Anna dem Jungen dürfen}.}

\begin{figure}[!htbp]
  \centering
  %\resizebox{\textwidth}{!}{
  \oneline{
    \begin{tabular}{cp{0.1em}cp{0.1em}cp{0.1em}cp{0.1em}c}
      VF && LSK && MF && RSK && NF \\
      \cmidrule{1-1}\cmidrule{3-3}\cmidrule{5-5}\cmidrule{7-7}\cmidrule{9-9}
      Schenken && will && Anna dem Jungen einen Ball. && -- && -- \\
      Einen Ball schenken && will && Anna dem Jungen. && -- && -- \\
      Dem Jungen einen Ball geschenkt && hat && Anna. && -- && -- \\
        * Anna schenken && will && dem Jungen einen Ball. && -- && -- \\
    \end{tabular}
  }
  \caption{Infiniter Verbteil (mit Ergänzungen) im Vorfeld}
  \label{fig:FeldermodelV2mitINFimVF}
\end{figure}

Wenn das Verb einteilig ist, aber dennoch durch die Stellung im\is{Vorfeld} VF hervorgehoben werden soll, dann können wir ein zweiteiliges Verb aus einer finiten Form von \textit{tun} und dem Infinitiv des eigentlich einteiligen Verbs bilden. Diese Struktur bezeichnet man als \textit{tun}-Periphrase. Die Aufteilung im Feldermodell wird in der Abbildung~\ref{fig:FeldermodelV2mitINFTUNimVF} dargestellt. Das Nachfeld bleibt der Übersichtlichkeit halber unbesetzt.

\begin{figure}[!htbp]
  \centering
  %\resizebox{\textwidth}{!}{
  \oneline{
    \begin{tabular}{cp{0.1em}cp{0.1em}cp{0.1em}cp{0.1em}c}
      VF && LSK && MF && RSK && NF \\
      \cmidrule{1-1}\cmidrule{3-3}\cmidrule{5-5}\cmidrule{7-7}\cmidrule{9-9}
      Schenken && tut && Anna dem Jungen einen Ball. && -- && -- \\
    \end{tabular}
  }
  \caption{Infiniter Verbteil im Vorfeld bei der \textit{tun}-Periphrase}
  \label{fig:FeldermodelV2mitINFTUNimVF}
\end{figure}

Solche Vorfeldbesetzungen bezeichnet \citet[202]{Zifonun2021DtEur} als selten, aber für das Deutsche typisch. Motiviert seien sie durch eine im Kontext aufgebaute Erwartungshaltung bezüglich des Verbs.\footnote{Als Beispiel führt \citet[202]{Zifonun2021DtEur} den folgenden Zeitungsbeleg auf, in dem es um eine Suche nach Rauschgift geht: \textit{Auch Drogenhunde wurden dabei eingesetzt. Gefunden wurde nichts.} Der Einsatz der Hunde bei der Suche nach Rauschgift baut die Erwartung eines Findens auf. Diese Erwartung motiviert die Voranstellung von \textit{gefunden}.}   

Das\is{Vorfeld} VF hat verschiedene Funktionen bei der Strukturierung von Informationen im Satz. Auf diese wird hier nur kurz eingegangen. Wir verweisen auf \citet{Musan2017} und mit Bezug auf den Unterricht in der Sekundarstufe auf \citet{Averintseva-KlischPeschel2014}. 

Bezieht sich ein Satz auf zuvor Gesagtes -- Sätze sind selten isoliert --, so nehmen wir im\is{Vorfeld} VF typischerweise die bekannte Entität wieder auf. Diese bekannte oder als bekannt vorausgesetzte Entität nennt man auch \textit{Thema} oder\is{Topik} \textit{Topik}. Das Satztopik ist diejenige Entität, über die etwas Neues ausgesagt wird bzw. von der die Äußerung primär handelt. Deutlich wird dies in Kontexten, in denen Sätze aufeinander folgen (\ref{ea:VFTopic}). 

\ea \label{ea:VFTopic} Gestern habe ich Paul\textsubscript{\textsc{neu}} getroffen. Den\textsubscript{\textsc{topik}} habe ich schon lange nicht mehr gesehen. 
\z 

Insbesondere in der Mündlichkeit kann das\is{Vorfeld} VF regulär unbesetzt bleiben, wenn das, worum es geht, also das Topik, in der Gesprächssituation präsent ist (\ref{ea:Topikdrop}). Hierauf gehen wir in Abschnitt~\ref{subsubsection:TestNOT} unter dem Stichwort\is{situative Ellipse} \textit{situative Ellipse} ein.\footnote{Bei Interesse an solchen Auslassungen empfehlen wir \citet{Fries1988}. Eine knappe Übersicht bieten \citet[413--418]{ZifHoffStr1997}.}

\ea \label{ea:Topikdrop} Wo ist Anna? -- \sout{Das} Weiß keiner.
\z 

Häufig werden temporale oder lokale Angaben vorangestellt, so wie \textit{gestern} im ersten Satz in (\ref{ea:VFTopic}). Diese stellen den örtlichen oder zeitlichen Rahmen des im Satz ausgedrückten Szenarios dar.

Zur besonderen Hervorhebung (\textit{Emphase}) können wir auch die neue Information ins\is{Vorfeld} VF stellen. Dies ist der  Fall bei \textit{ungemütlich} im zweiten Satz in (\ref{ea:VFNichtTopik}). 

\ea \label{ea:VFNichtTopik} Paul ist ein freundlicher Mensch. Ungemütlich\textsubscript{\textsc{neu}} wird er\textsubscript{\textsc{topik}} nur, wenn er Hunger hat.
\z 

Aus informationsstrukturellen Gründen kann im Vorfeld auch ein semantisch leeres \textit{es} als sogenannter\is{es@\textit{es}!Platzhalter} \textit{Platzhalter} (auch \textit{Expletivum}) stehen. Es hält den im Verbzweitsatz obligatorisch zu besetzenden Platz im Vorfeld besetzt. Dieses \textit{es} ist wie in (\ref{ea:Platzhalter_es}) keine Ergänzung des Verbs, keine Angabe und kein infiniter Prädikatsteil.   

\ea \label{ea:Platzhalter_es} Anna hatte gestern Geburtstag. \textit{Es} gratulierten ihr alle ihre Freunde.
\z 

Dass\is{es@\textit{es}!Platzhalter} \textit{es} kein Subjekt ist, zeigt sich insbesondere an der finiten Verbform \textit{gratulierten} (3. Pers./Pl.). Sie kongruiert nicht mit \textit{es} (3. Pers./Sgl.). Es handelt sich um ein reines Stellungsglied. Dieses erlaubt es, eine Angabe oder typischerweise die erste Verbergänzung, die im Vorfeld stehen würde, ins Mittelfeld zu stellen. Je später eine Einheit genannt wird, umso mehr steht sie im Zentrum der Aufmerksamkeit (\textit{Fokus}). 

Formal erkennt man\is{es@\textit{es}!Platzhalter} Platzhalter daran, dass sie nicht vom jeweiligen Verb gefordert werden und somit aus dem Satz verschwinden, sobald eine Ergänzung oder Angabe das Vorfeld besetzt (\ref{ea:Wegfall_Platzhalter_es}). 

\ea \label{ea:Wegfall_Platzhalter_es} Anna hatte gestern Geburtstag. Ihr gratulierten alle ihre Freunde.
\z 

Wenn \textit{es} hingegen formal als\is{es@\textit{es}!formales Subjekt} erste Ergänzung wie von den Witterungsverben regiert ist, dann steht es bei der Voranstellung einer anderen Einheit wie in (\ref{ex:regnete_es}) im Mittelfeld.\footnote{Wenn \textit{es} formal als zweite Ergänzung regiert wird wie in den festen Wendungen \textit{es zu etwas bringen} oder \textit{es sich leicht/schwer machen}, dann kann \textit{es} ohnehin nicht im Vorfeld stehen. Eine kurze Übersicht zum Status von \textit{es} bieten \citet[239--244]{HelbigBuscha}. Auf \textit{es} als Platzhalter kommen wir noch einmal beim sogenannten unpersönlichen Passiv im Abschnitt~\ref{subsection:Passiv} zurück.}

\ea \label{ea:form_Subjekt}
\ea \label{ea:es_regnete} \textit{Es} regnete gestern den ganzen Tag.
\ex \label{ex:regnete_es} Gestern regnete \textit{es} den ganzen Tag.
\z 
\z 

Noch vor dem VF können jedoch sogenannte Konjunktionen wie \textit{und}, \textit{oder}, \textit{aber} oder \textit{denn} stehen. Eine Übersicht zu Konjunktionen als Wortart geben wir in Abschnitt~\ref{subsubsection:Konjunktion}. Ihre Funktion ist der inhaltliche Anschluss eines Satzes an einen vorhergehenden Satz. Wie in den Beispielen in (\ref{ea:Anschlussstelle}) ersichtlich wird, stehen Konjunktionen in dieser Funktion vor dem Vorfeld.  

\ea \label{ea:Anschlussstelle}
\ea \label{ea:ANund} Es regnet. \textit{Aber} Anna geht spazieren.
\ex \label{ex:denn} Es muss Nachtfrost gegeben haben. \textit{Denn} alle Pfützen sind zugefroren.
\z 
\z 

Diese Konjunktionen besetzen ein Feld vor dem Vorfeld. Da es um die Anschlussfunktion des Satzes geht, bezeichnen wir das Feld mit \citet[65]{Woellstein2014} als\is{Anschlussfeld} \textit{Anschlussfeld} (AF).\footnote{\citet[80]{Agel2017} schlägt den Begriff \textit{Zwischenstelle} vor. \citet[398]{Schaefer2018_3Aufl} schlägt \textit{Konnektorfeld} vor.} Das erweiterte Feldermodell wird in \ref{fig:FeldermodelmitAF} anhand der Sätze in (\ref{ea:ANund}) und (\ref{ex:denn}) dargestellt. Zur Illustration besetzen wir im ersten Satz auch das NF. 

\begin{figure}[!htbp]
  \centering
   %\resizebox{\textwidth}{!}{
   \oneline{
  \begin{tabular}{cp{0.1em}cp{0.1em}cp{0.1em}cp{0.1em}cp{0.1em}c}
    AF && VF && LSK && MF && RSK && NF \\
    \cmidrule{1-1}\cmidrule{3-3}\cmidrule{5-5}\cmidrule{7-7}\cmidrule{9-9}\cmidrule{11-11}
    Aber && Anna && geht && -- && spazieren && trotz des Regens. \\
    Denn && überall && liegt && Raureif. && -- && -- \\
  \end{tabular}
  }
  \caption{Anschlussfeld vor dem Vorfeld}
  \label{fig:FeldermodelmitAF}
\end{figure}

Bevor wir uns in Abschnitt~\ref{subsection:Mittelfeld} mit der Besetzung des Mittelfelds beschäftigen, gehen wir in dem folgenden Exkurs~\ref{Exkurs:Mehrfache VFB} auf Stellungen ein, bei denen nach dem Anschlussfeld und vor dem Vorfeld weitere Konstituenten stehen. Solche Strukturen sehen nach einer mehrfachen Vorfeldbesetzung und damit nach einer Abweichung von der Regel aus, dass im Vorfeld \textit{eine} Satzkonstituente steht.


\begin{Exkurs}{Mehrfache Vorfeldbesetzung?}\label{Exkurs:Mehrfache VFB}
    
Zwischen dem AF und dem VF kann man eine weitere Konstituente platzieren. Das illustriert das von \citet[70]{Hoehle19862018} entnommenene Beispiel in  (\ref{ea:Linksversetzung}). 

\ea \label{ea:Linksversetzung} Seinen Hund, den sollte man anständig behandeln.
\z 

Es handelt sich um eine sogenannte\is{Linksversetzung} \textit{Linksversetzung}, für die man ein \textit{Vor-Vorfeld} vor dem Vorfeld annimmt. In (\ref{ea:Linksversetzung}) steht \textit{seinen Hund} in diesem Vor-Vorfeld. Die linksversetzte Einheit wird durch einen \textit{d}-Ausdruck wieder aufgenommen. In (\ref{ea:Linksversetzung}) steht \textit{den} im Vorfeld und nimmt den Ausdruck im Vor-Vorfeld wieder auf. Der linksversetzte Ausdruck \textit{sein-en Hund} kongruiert mit dem ihn wieder aufnehmenden Ausdruck \textit{d-en}. Solche Linksversetzungen dienen der Festlegung eines\is{Topik} Satztopiks, d.\,h. eines Gegenstands, um den es in dem entsprechenden Satz geht. Dieser Gegenstand ist im jeweiligen Diskurskontext bereits eingeführt.\footnote{Zur Vertiefung empfehlen wir \citet[54--56]{Woellstein2014} sowie \citet[142--153]{AltmannHofman}.} 

Im Gegensatz zu solchen kongruenten Linksversetzungen gibt es auch nicht-kongruierende Konstruktionen mit einem linksversetzten Ausdruck im Nominativ wie \textit{der Hund} in (\ref{ea:hängendesTopik}). 

\ea \label{ea:hängendesTopik} Der Hund, den sollte man anständig behandeln.
\z 

Diese Art der Linksversetzung bezeichnet man auch als \textit{hängende Topikkonstruktion} (im Englischen \textit{hanging topic}). Sie dient der Einführung eines ganz neuen Diskurstopiks, d.\,h. eines Gegenstands, um den es in dem jeweiligen Kontext noch gar nicht ging. Der linksversetzte Ausdruck, das hängende Topik, ist durch fehlende Kongruenz und einen intonatorischen Bruch weniger integriert als bei der Linksversetzung in (\ref{ea:Linksversetzung}). Daher schlägt \citet[71]{Woellstein2014} vor, für das hängende Topik ein \textit{Außenfeld} vor dem Vor-Vorfeld anzunehmen. Bei \citet[1580]{ZifHoffStr1997} besetzen sowohl das hängende Topik als auch die Linksversetzung die gleiche Position im linken Außenfeld, das allerdings unter anderem auch eine Position für die koordinierenden Ausdrücke umfasst. Wir lassen offen, ob man dem Unterschied zwischen der Linksversetzung und dem hängenden Topik durch die Annahme unterschiedlicher Felder (Außenfeld vs. Vor-Vorfeld) besser gerecht wird. Wesentlich ist, dass es sich weder bei der Linksversetzung noch beim hängenden Topik um doppelte Vorfeldbesetzungen handelt.   

Es gibt weitere Fälle wie in (\ref{ea:Verbdritt?}), in denen das Verb in der LSK an dritter Stelle steht. Zur Veranschaulichung klammern wir die beiden Wortgruppen vor dem finiten Verb in der LSK ein.

\ea\label{ea:Verbdritt?}
\ea \label{ea:LOCXV} [In anderen Städten] [das] gibt es nicht.\footnote{Der Beleg ist \citet[2]{Schalowski2017} entnommen und stammt aus einem Fernsehbericht ohne multietnolektalen Hintergrund.}
\ex \label{ex:TEMPXV} [Ab jetzt] [ich] kriege immer zwanzig Euro.\footnote{Der Beleg ist \citet[84]{Freywald_al_2015} entnommen und stammt aus dem KiezDeutsch-Korpus.}
\z 
\z 

Stellungen wie in (\ref{ea:Verbdritt?}) finden sich vor allem in der gesprochenen Sprache, und zwar sowohl in sogenannten multiethnolektalen Sprechweisen wie dem Kiezdeutsch als auch bei monolingualen Sprecherinnen und Sprechern des Deutschen. Zumindest in der geschriebenen Standardsprache wäre die Verbzweitstellung zu erwarten: Also (\ref{ea:LOCV}) statt (\ref{ea:LOCXV}) oder eine Linksversetzung ins Vor-Vorfeld mit anschließender Wiederaufnahme durch \textit{da} im Vorfeld in (\ref{ex:LOCdaV}) und (\ref{ex:TEMPV}) statt (\ref{ex:TEMPXV}) oder eine Linksversetzung ins Vor-Vorfeld mit anschließender Wiederaufnahme durch \textit{da} im Vorfeld wie in (\ref{ex:TEMPdaV}). 

\ea\label{ea:VerbdrittzuVerbzweit}
\ea \label{ea:LOCV} [In anderen Städten] gibt es [das] nicht.
\ex \label{ex:LOCdaV} [In anderen Städten], [da] gibt es [das] nicht.
\ex \label{ex:TEMPV} [Ab jetzt] kriege [ich] immer zwanzig Euro.
\ex \label{ex:TEMPdaV} [Ab jetzt], [da] kriege [ich] immer zwanzig Euro.
\z 
\z 

Aber kommen wir wieder auf die Belege in (\ref{ea:Verbdritt?}) zurück. Mit der ganz links platzierten Einheit, in der Regel Angaben zu Raum und Zeit, setzen Sprecherinnen und Sprecher ein sogenanntes Rahmentopik, das den Gültigkeitsbereich – im Gegensatz zu anderen möglichen Gültigkeitsbereichen – der Satzaussage einschränkt. Die Aussage in (\ref{ea:LOCXV}) gilt nur \textit{in anderen Städten} und nicht etwa in der eigenen Stadt. Die Aussage in (\ref{ex:TEMPXV}) gilt \textit{ab jetzt} und nicht etwa ab morgen. Wir kommen auf solche satzbezogenen Adverbiale in Abschnitt~\ref{subsection:Adverbial} zurück. 

In den Sätzen in (\ref{ea:Verbdritt?}) stehen die rahmensetzenden Topiks vor den Satztopiks \textit{das} in (\ref{ea:LOCXV}) und \textit{ich} in (\ref{ex:TEMPXV}). Danach wird ab dem finiten Verb in der LSK die neue Information über das Satztopik gegeben, die für den durch den Rahmensetzer ausgedrückten Gültigkeitsbereich gilt. Das Ergebnis solcher Informationsstrukturierungen ist die sogenannte Verbdrittstellung, da sowohl das Rahmentopik als auch das Satztopik vor dem finiten Verb, eben in der Topikposition stehen.\footnote{Cf. \citet[19--24]{Schalowski2017} und \citet[88--91]{Freywald_al_2015}, die dasselbe Phänomen auch in anderen germanischen Verbzweit-Sprachen wie Norwegisch, Schwedisch und Niederländisch betrachten.} Solche Stellungen finden sich bereits in älteren Sprachstufen des Deutschen.\footnote{\citet[235]{Demske_Wiese2016} vermuten daher, dass aktuelle Verbdrittstellungen wie in (\ref{ea:Verbdritt?}) nicht das Ergebnis neuerer Entwicklungen sein müssen, sondern optionale Reste aus Sprachstufen, in denen die Verbzweitstellung noch nicht etabliert war. Dagegen hält  \citet[35--37]{Breitbarth2023} aufgrund ihrer Datenauswertung solche Verbdrittsätze für das Ergebnis einer neueren Entwicklung.} Heute werden sie nach \citet[167]{Freywald_al_2023} \gqq{oft mit multiethnolektalen Sprechweisen wie Kiezdeutsch in Verbindung gebracht}, wo sie zwar eine besondere Dynamik entfalten, jedoch nicht auf diese beschränkt sind. Während \citet[166]{Freywald_al_2023}\footnote{Sie berufen sich auf Arbeiten von \citet{Walkden2017} und \citet{Bunk2020}, die jedoch auf anderen grammatischen Modellen beruhen.} davon ausgehen, dass sowohl das Rahmentopik als auch das Satztopik das Vorfeld besetzen, beschreiben wir solche Stellungen analog zu der bereits erwähnten hängenden Topikkonstruktion mit einem Außenfeld. Anstelle eines neuen Diskurstopiks besetzt das Rahmentopik eine Stelle vor dem Vorfeld.\footnote{Hierfür spricht besonders die Prosodie solcher Verbdrittstellungen, die \citet{Breitbarth2022} untersucht. Sie zeigt, dass die Rahmentopiks prosodische Einheiten bilden, die nicht in den folgenden Verbzweit-Satz integriert sind.} 

Die in (\ref{ea:Verbdritt?}) aufgeführten Belege sind zu unterscheiden von Stellungen wie in (\ref{ea:scheinbareMVB}).

\ea \label{ea:scheinbareMVB} [Der Universität] [zum Geburtstag] gratulierte auch Bundesminister Dorothee Wilms [\ldots]\footnote{Der Beleg ist \citet[87]{Duerscheid1989} entnommen.}  
\z 

In Abschnitt~\ref{section:Voranstellungstest} kommen wir auf solche scheinbaren mehrfachen Vorfeldbesetzungen zurück und werden mit \citet{Mueller2005} dafür argumentieren, die beiden Einheiten \textit{der Universität} und \textit{zum Geburtstag} zu \textit{einer} größeren Einheit zusammenzufassen analog zu (\ref{ea:VK_VF}).

\ea\label{ea:VK_VF} [Der Universität zum Geburtstag gratuliert] hat auch Dorothee Wilms.
\z  

In (\ref{ea:VK_VF}) stehen die beiden Ergänzungen des Verbs \textit{der Universität} und \textit{zum Geburtstag} zusammen mit der sie regierenden infiniten Verbform \textit{gratuliert} im Vorfeld und bilden \textit{eine} Einheit, eben eine Verbgruppe. Diese steht im Vorfeld. Wir haben bereits darauf hingewiesen, dass es für solche Wortgruppen in der traditionellen Schulgrammatik keinen etablierten Begriff gibt und kommen in Abschnitt~\ref{section:VP} unter dem Stichwort \textit{infinite Verbalphrase} hierauf zurück. 


\end{Exkurs}

\subsection{Mittelfeld}\label{subsection:Mittelfeld}

Das\is{Mittelfeld} MF von Verbzweitsätzen wird durch die LSK geöffnet. Wenn eine RSK vorhanden ist, so schließt sie das MF. Im MF können beliebig viele Konstituenten stehen. Die Stellung der Konstituenten ist im MF relativ frei. Da es uns nur um die Stellungsmöglichkeiten im MF geht, verzichten wir in der Darstellung auf die Besetzung des AF und des NF. In der Abbildung~\ref{fig:Mittelfeldstellungen} zeigen wir die Umstellungsmöglichkeiten der Konstituenten \textit{Anna}, \textit{dem Jungen} und \textit{den Ball} im Mittelfeld.


\begin{figure}[!htbp]
  \centering
   %\resizebox{\textwidth}{!}{
   \oneline{
  \begin{tabular}{cp{0.1em}cp{0.1em}cp{0.1em}c}
    VF && LSK && MF && RSK \\
    \cmidrule{1-1}\cmidrule{3-3}\cmidrule{5-5}\cmidrule{7-7}
    Gestern && hat && Anna dem Jungen den Ball && geschenkt. \\
    Gestern && hat && dem Jungen Anna den Ball && geschenkt. \\
   Gestern && hat && dem Jungen den Ball Anna && geschenkt. \\
   Gestern && hat && Anna den Ball dem Jungen && geschenkt. \\
    Gestern && hat && den Ball Anna dem Jungen && geschenkt. \\
     Gestern && hat && den Ball dem Jungen Anna && geschenkt. \\
  \end{tabular}
  }
  \caption{Stellungsmöglichkeiten der Konstituenten im Mittelfeld}
  \label{fig:Mittelfeldstellungen}
\end{figure}

Die Abfolge im ersten Satz ist geläufiger als die in den anderen Sätzen. Möglich und korrekt, wenn auch mit besonderer Betonung und in bestimmten, eingeschränkten Kontexten, sind jedoch auch alle übrigen Abfolgen. 

Die\is{Mittelfeld!Reihenfolge} verschiedenen grammatischen, semantischen, prosodischen und pragmatischen Faktoren, die für die Abfolge im MF verantwortlich sind, werden wir hier nicht weiter darstellen.\footnote{Zur theoretischen Auseinandersetzung mit der Gewichtung und Hierarchie der Faktoren verweisen wir auf \citet{Hoehle19822018} (erstmals schon 1982 veröffentlicht) und auf \citet{Reis1987}.} Bevor es in Abschnitt~\ref{subsection:Nachfeld} um die Besetzung des Nachfelds geht, stellen wir in dem folgenden Exkurs~\ref{Exkurs:Stellung_MF} einige Stellungspräferenzen im Mittelfeld vor. Denn diese können insbesondere für den Fremd- und Zweitsprachenunterricht hilfreich sein. Auf diese Relevanz macht auch \citet[143]{Tahiri2022} mit dem folgenden Beleg eines marokkanischen DaF-Lerners aufmerksam: \textit{Der Dieb lief in Richtung des Nagels mit großen Schritten und dann trat er auf den Nagel mit großer Wucht.} 


\begin{Exkurs}{Stellungspräferenzen im Mittelfeld}\label{Exkurs:Stellung_MF}

\largerpage
Für die Abfolge\is{Mittelfeld!Reihenfolge} im Mittelfeld gelten folgende Tendenzen: Nominativ-NP vor Dativ-NP vor Akkusativ-NP vor PP (\ref{ea:Nom_Dat_Akk_PP}), pronominale vor substantivischer NP (\ref{ex:Prn_NP}) und definite NP vor indefiniter NP (\ref{ex:def_indef}). 

\ea \label{ea:Abfolge_MF}
\ea \label{ea:Nom_Dat_Akk_PP} Gestern hat [die Oma]\textsubscript{\textsc{nominativ-np}} [ihrem Enkel]\textsubscript{\textsc{dativ-np}} [den Kakao]\textsubscript{\textsc{akkusativ-np}} [ans Bett]\textsubscript{\textsc{pp}} gebracht.
\ex \label{ex:Prn_NP} Gestern hat [ihm]\textsubscript{\textsc{pronominal}} [die Oma]\textsubscript{\textsc{substantivisch}} gratuliert.
\ex \label{ex:def_indef} Gestern hat [dem Enkel]\textsubscript{\textsc{definit}} auch [ein Nachbar]\textsubscript{\textsc{indefinit}} gratuliert.
\z 
\z



Informationen, die als bekannt vorausgesetzt werden oder vorerwähnt sind, werden definit oder durch ein Pronomen ausgedrückt. Sie\is{Mittelfeld!Reihenfolge} stehen vor neu eingeführten und häufig betonten Informationen, die indefinit ausgedrückt werden.  

Die Abfolge\is{Mittelfeld!Reihenfolge} substantivischer NPs in (\ref{ea:Nom_Dat_Akk_PP}) illustriert die folgende Normalabfolge: Nominativ-NP vor Dativ-NP vor Akkusativ-NP. Eine mögliche Erklärung sind die in den Abschnitten~\ref{subsubsection:TestINSP} und \ref{subsection:Verbklassen} besprochenen semantischen Rollen, insbesondere das Agensgefälle, das mit Belebtheit einhergeht. Bei der Mehrzahl der dreiwertigen Verben realisiert die Nominativ-NP ein Agens (Handlungs- oder Tätigkeitsträger). Die NP im Akkusativ realisiert das Patiens (Handlungsgegenstand). Die Dativ-NP kodiert einen Betroffenen oder Adressaten. Die Rolle des Betroffenen oder Adressaten ist agentiver als die des Patiens und steht deshalb vor diesem.

Ein Indiz für die Gewöhnlichkeit der\is{Mittelfeld!Reihenfolge} Abfolge Dativ-NP vor Akkusativ-NP ist die Tatsache, dass die NP im Akkusativ (\ref{ea:Dat_AKK}) und die NP im Dativ (\ref{ex:DAT_Akk}) den Satzakzent und damit Informationsgewicht tragen können. Wenn die Akkusativ-NP der Dativ-NP (\ref{ex:Akk_DAT}) vorausgeht, kann die Akkusativ-NP nicht (ohne speziellen Kontrast) den Satzakzent tragen (\ref{ex:AKK_Dat}). Die Betonung wird durch Großbuchstaben wiedergegeben.  

\ea \label{ea:Dat-NPvorAkk-NP}
\ea [] {Gestern hat Anna dem Jungen den BALL geschenkt.}\label{ea:Dat_AKK}
\ex [] {Gestern hat Anna dem JUNGEN den Ball geschenkt.}\label{ex:DAT_Akk}
\ex [] {Gestern hat Anna den Ball dem JUNGEN geschenkt.}\label{ex:Akk_DAT}
\ex [*] {Gestern hat Anna den BALL dem Jungen geschenkt.}\label{ex:AKK_Dat}
\z 
\z 

Bei einer relativ kleinen Gruppe von Verben\is{Mittelfeld!Reihenfolge} steht jedoch die Akkusativ-NP gewöhnlich vor der Dativ-NP. Dies betrifft nach \citet[134--135]{Wegener1995} Nominalisierungsverbgefüge (vgl. Abschnitt \ref{subsection:IdiomeNVGFVG}) wie in (\ref{ea:FVG}) und Handlungsverben, mit denen ausgedrückt wird, dass eine Größe einer anderen gleich oder eben nicht gleichgestellt wird. Hierzu gehören \zB \textit{angleichen}, \textit{unterwerfen} und \textit{vorziehen} in (\ref{ex:vorziehenAkkDat}).

\ea \label{ea:Akk-NPvorDat-NP}
\ea \label{ea:FVG} Gestern hat die Lehrerin den Schüler einer Prüfung unterzogen.
\ex \label{ex:vorziehenAkkDat} Immer zieht die Oma ihren Enkel ihrem Sohn vor.
\z 
\z 

Die\is{Mittelfeld!Reihenfolge} Normalabfolge Dativ-NP vor Akkusativ-NP können wir umkehren, wenn das informative Hauptgewicht, der \textit{Fokus}, auf der Dativ-NP liegt. So kann die Akkusativ-NP vor der Dativ-NP stehen, wenn letztere erfragt wird (\ref{ea:MFAkkDat}). Sie trägt dann auch den Hauptakzent. Dies entspricht dem zweiten \textit{Behaghelschen Gesetz}, demzufolge \gqq{das Wichtige später steht als das Unwichtige}. Denn nach \citet[4]{Behaghel1932SyntaxStellung} ist das Wichtige das, \gqq{was zuletzt noch im Ohr klingen soll}, also möglichst weit hinten steht. 

\ea \label{ea:MFAkkDat} WEM hat Anna den Ball geschenkt? -- Anna hat den Ball dem JUNGEN geschenkt.
\z 

Eine\is{Mittelfeld!Reihenfolge} Umkehrung ist auch dann üblich, wenn die Akkusativ-NP definit ist und die Dativ-NP indefinit (\ref{ea:DefIndef}). Der Hintergrund ist, wie oben bereits erwähnt, dass die neue (indefinit eingeführte) Information in der Regel nach der schon bekannten (definit eingeführten) Information steht.

\ea \label{ea:DefIndef} Anna hat den Ball\textsubscript{\textsc{definit}} einem Jungen\textsubscript{\textsc{indefinit}} geschenkt.
\z 

Wird die Akkusativ-NP als Pronomen realisiert, dann steht das Akkusativ-Pronomen vor der Dativ-NP\is{Mittelfeld!Reihenfolge}. Dies ist unabhängig davon, ob die Dativ-NP auch als Pronomen ausgedrückt wird (\ref{ea:MFAkkPrnDat}) oder nicht (\ref{exMFAkkPrnDatNP}). 

\ea 
\ea \label{ea:MFAkkPrnDat} (Das ist Pauls Ball.) Anna hat ihn ihm geschenkt.
\ex \label{exMFAkkPrnDatNP} (Das ist Pauls Ball.) Anna hat ihn dem Jungen geschenkt.
\z
\z 

Bei der Abfolge Akkusativ-Pronomen vor Dativ-NP\is{Mittelfeld!Reihenfolge} steht die bekannte (und deshalb als Pronomen realisierte) Information wieder vor der neuen Information. Diese Stellung folgt auch der Stellungspräferenz kurz vor lang. Die typische Abfolge kurz vor lang bezeichnet \citet[6]{Behaghel1932SyntaxStellung} als \textit{Gesetz der wachsenden Glieder}.\footnote{\citet[137]{Wegener1995} erklärt die umgekehrte Reihenfolge im Pronominalbereich grundsätzlich durch dieses Behaghelsche Gesetz. Einsilbige Pronomen stehen vor zweisilbigen: \textit{Sie hat sie ihnen gezeigt}. Kurze Vokale stehen vor langen: \textit{Sie hat mich dir vorgestellt}. Auslautende Vokale stehen vor auslautenden Konsonanten: \textit{Er hat sie ihm erklärt}.} 

Unabhängig von der Informationsgewichtung steht im MF eine Akkusativ-NP vor einer Ergänzung, die durch eine PP ausgedrückt wird (\ref{ea:MFAkkPP}).

\ea \label{ea:MFAkkPP} Anna hat den Jungen zur Schule gebracht.
\z 

Temporale Angaben\is{Mittelfeld!Reihenfolge} stehen vor einer Dativ- oder Akkusativ-NP (\ref{ea:TempDat}). Lokale Angaben folgen einer Dativ- oder Akkusativ-NP (\ref{ex:AkkLok}). Einer PP-Ergänzung gehen sie voraus (\ref{ex:LokPobj}).

\ea 
\ea \label{ea:TempDat} Anna hat gestern dem Jungen geholfen/gestern den Jungen gesehen.
\ex \label{ex:AkkLok} Anna hat den Ball auf dem Spielplatz gefunden.
\ex \label{ex:LokPobj} Anna wartet auf dem Spielplatz auf Paul.
\z 
\z 

Unter den Angaben gibt es im MF die\is{Mittelfeld!Reihenfolge} Abfolgetendenz temporal vor kausal vor modal vor lokal (\ref{ea:TeKaMoLo}). Diese Abfolge bezeichnet man im fremdsprachlichen Unterricht auch als \textit{TeKaMoLo}-Prinzip: \textit{Te} steht für temporal, \textit{Ka} für kausal, \textit{Mo} für modal und \textit{Lo} für lokal.

\ea \label{ea:TeKaMoLo} Anna hat gestern\textsubscript{\textsc{te}} trotz des Regens\textsubscript{\textsc{ka}} lange\textsubscript{\textsc{mo}} auf dem Spielplatz\textsubscript{\textsc{lo}} gewartet.
\z 

Die hier dargestellten Abfolgen im MF gelten nicht nur für Verbzweitsätze, sondern ebenso für die MF von Verberst- und Verbletztsätzen.

\end{Exkurs}


\subsection{Nachfeld}\label{subsection:Nachfeld}

Die Besetzung des\is{Nachfeld} NF ist fakultativ. Sind unmittelbare Satzkonstituenten jedoch als Nebensätze oder Infinitivgruppen realisiert, so stehen sie typischerweise im NF. Auch durch PP realisierte Satzkonstituenten können im Nachfeld stehen. Diese Besetzungsmöglichkeiten stellen wir in \ref{fig:Nachfeld} dar.


\begin{figure}[!htbp]
  \centering
  %\resizebox{\textwidth}{!}{
  \oneline{
    \begin{tabular}{cp{0.1em}cp{0.1em}cp{0.1em}cp{0.1em}c}
      VF && LSK && MF && RSK && NF \\
      \cmidrule{1-1}\cmidrule{3-3}\cmidrule{5-5}\cmidrule{7-7}\cmidrule{9-9}
      Anna && hat && Paul einen Ball && geschenkt, && weil er Geburtstag hat. \\
      Anna && hat && Paul && gebeten, && ihr zu helfen.\\ 
      Anna && weist && immer && hin && auf Pünktlichkeit.
    \end{tabular}
  }
  \caption{Besetzung des Nachfeldes}
  \label{fig:Nachfeld}
\end{figure}

Wenn die RSK wie in (\ref{ea:MFoderNF?}) nicht besetzt ist, dann ist auf den ersten Blick nicht klar, ob die Konstituente \textit{er will an die Ostsee fahren} im MF oder im NF steht.

\ea \label{ea:MFoderNF?} Paul sagt, er will an die Ostsee fahren.
\z 

Hier hilft es, aus dem einteiligen ein zweiteiliges Verb zu bilden, \zB durch die analytische Perfektbildung (\textit{hat gesagt}). Die Stellung von \textit{er will an die Ostsee fahren} im Nachfeld (\ref{ea:NeSaimNF}) scheint viel akzeptabler als die im Mittelfeld (\ref{ex:*NeSaimMF}).  

\newpage
\ea \label{ea:TestMFoderNF}
\ea [] {Paul hat gesagt, er will an die Ostsee fahren.}\label{ea:NeSaimNF}
\ex [?] {Paul hat, er will an die Ostsee fahren, gesagt.}\label{ex:*NeSaimMF}
\z 
\z 

Die fragliche Akzeptabilität von (\ref{ex:*NeSaimMF}) spricht dafür, den Ergänzungssatz \textit{er will an die Ostsee fahren} im NF zu verorten. 

Bestimmte Satzkonstituenten werden statt dem Mittelfeld im Nachfeld positioniert, was man als\is{Nachfeld!Ausklammerung} \textit{Ausklammerung} bezeichnet. So werden Relativsätze, die Teil (Attribut) einer Satzkonstituente im MF sind (\ref{ea:RelMF}), normalerweise ins NF ausgelagert (\ref{ex:RelNF}). Dadurch wird die Satzklammer nicht überdehnt und die Sprachverarbeitung erleichtert. 

\ea \label{ea:AuslagerungReklativsatz}
\ea \label{ea:RelMF} Anna hat gestern Paul, den sie schon lange nicht mehr gesehen hat, getroffen.
\ex \label{ex:RelNF} Anna hat gestern Paul getroffen, den sie schon lange nicht mehr gesehen hat.
\z 
\z 

Insbesondere PPs können gut\is{Nachfeld!Ausklammerung} ausgeklammert werden. Häufig, aber nicht immer, ist dies wie in (\ref{ea:AUsklammerungPP}) durch ein besonderes Informationsgewicht motiviert. 

\ea \label{ea:AUsklammerungPP} Paul hat allen gedankt für die tollen Geschenke.
\z

Auch hier gilt, dass der wichtige Teil der Information möglichst am Ende steht und somit, nach \citet[4]{Behaghel1932SyntaxStellung}, \gqq{zuletzt noch im Ohr klingen soll}. 

Da unter verschiedenen stilistischen und informationsstrukturellen Voraussetzungen tendenziell alle Satzkonstituenten das NF besetzen können, führen wir hier die wichtigsten Einheiten auf, die das NF nicht besetzen können. Es handelt sich um Abtönungspartikeln (vgl. Abschnitt~\ref{subsection:Partikeln}), die Satznegation \textit{nicht} und Reflexivpronomen. Dies wird in Abbildung~\ref{fig:unmöglicheNachfeldbesetzung} dargestellt. 

\begin{figure}[!htbp]
  \centering
  %\resizebox{\textwidth}{!}{
  \oneline{
    \begin{tabular}{cp{0.1em}cp{0.1em}cp{0.1em}cp{0.1em}c}
      VF && LSK && MF && RSK && NF \\
      \cmidrule{1-1}\cmidrule{3-3}\cmidrule{5-5}\cmidrule{7-7}\cmidrule{9-9}
      * Das && habe && ich dir && gesagt && doch \\
      * Paul && hat && -- && gelernt && nicht.\\ 
      * Anna && hat && -- && geirrt && sich.
    \end{tabular}
  }
  \caption{Unmögliche Besetzung des Nachfeldes}
  \label{fig:unmöglicheNachfeldbesetzung}
\end{figure}

In Abschnitt~\ref{subsection:Vorfeld} haben wir darauf hingewiesen, dass es zwischen dem AF und dem VF ein sogenanntes Vor-Vorfeld für Linksversetzungen gibt. Ebenso gibt es auch\is{Rechtsversetzung} Rechtsversetzungen und demzufolge nach dem NF ein sogenanntes \textit{rechtes Außenfeld}. Dessen Besetzung wird in dem von \citet[73]{Woellstein2014} übernommenen Satz (\ref{ea:rechtesAußenfeld}) illustriert.

\ea \label{ea:rechtesAußenfeld} Wir wollten doch feiern mit dir, mein Sohn, weil das dein 18. Geburtstag ist.
\z 

Das Pronomen der PP \textit{mit dir} im NF wird durch die NP \textit{mein Sohn} im rechten Außenfeld \textit{mein Sohn} als eine Art Ansprache wieder aufgenommen bzw. expliziert. Darauf folgt im sogenannten \textit{weiten Nachfeld} noch der Angabe-Satz \textit{weil das dein 18. Geburtstag} ist. Wir gehen nicht weiter auf Rechtsversetzungen ein und verweisen für eine Übersicht auf \citet[153--157]{AltmannHofman}.  


\section{Verberstsätze}\label{section:Verberstsatz}

Verberstsätze\is{Verberstsatz} sind dadurch gekennzeichnet, dass das Finitum an erster Stelle steht, weshalb man sie auch \textit{Stirn}sätze nennt. Mit Stirnsätzen drücken wir\is{Entscheidungsfragesatz} Entscheidungsfragen (\ref{ea:V1Entscheidungsfrage}) und Aufforderungen im Imperativ (\ref{ex:V1Imperativ}) aus.

\ea \label{ea:Verberst}
\ea \label{ea:V1Entscheidungsfrage} Hat Anna dem Jungen einen Ball geschenkt?
\ex \label{ex:V1Imperativ} Schenk dem Jungen einen Ball!
\z 
\z 

Mit\is{Entscheidungsfragesatz} Entscheidungsfragen wird erfragt, ob die entsprechende Aussage (Verbzweitsatz) wahr ist. Entscheidungsfragen nennt man in der Schule auch \textit{Ja/Nein-Fragen}, weil man nur mit \textit{ja} oder \textit{nein} auf sie antworten kann. Mit Imperativsätzen wird dem Gegenüber das Wahrwerden der entsprechenden Aussage aufgetragen. 

Topologisch zeichnen sich\is{Verberstsatz} Verberstsätze dadurch aus, dass sie kein VF haben. Bis auf die Abwesenheit eines VF sind die Felder im Verberstsatz wie im Verbzweitsatz. Das Finitum öffnet als linke Satzklammer das MF. Falls vorhanden, wird es durch den infiniten Verbrest als rechte Satzklammer geschlossen. Nach der RSK gibt es genauso wie bei Verbzweitsätzen ein besetzbares NF. Bei Imperativsätzen kann in Abhängigkeit von der Valenz des Verbs minimal die LSK besetzt sein. Dies wird in \ref{fig:Verberstfelder} dargestellt.

\begin{figure}[!htbp]
  \centering
  \oneline{
  \begin{tabular}{cp{0.1em}cp{0.1em}cp{0.1em}c}
    LSK && MF && RSK && NF \\
    \cmidrule{1-1}\cmidrule{3-3}\cmidrule{5-5}\cmidrule{7-7}
    Hat && Anna ihm einen Ball && geschenkt && zum Geburtstag? \\
    Schenk && ihm einen Ball! && -- && -- \\
    Schlaf! && -- && -- && -- \\ 
  \end{tabular}
  }
  \caption{Feldermodell für Verberstsätze}
  \label{fig:Verberstfelder}
\end{figure}

In Abschnitt~\ref{subsection:Vorfeld} haben wir auf das\is{Anschlussfeld} Anschlussfeld (AF) hingewiesen. Dieses gibt es auch bei Verberstsätzen und wie bei Verbzweitsätzen kann es von Konjunktionen besetzt werden, was in der Abbildung~\ref{fig:VerberstfeldermitAF} gezeigt wird.

\begin{figure}[!htbp]
  \centering
  %\resizebox{\textwidth}{!}{
  \oneline{
    \begin{tabular}{cp{0.1em}cp{0.1em}cp{0.1em}cp{0.1em}c}
      AF && LSK && MF && RSK && NF \\
      \cmidrule{1-1}\cmidrule{3-3}\cmidrule{5-5}\cmidrule{7-7}\cmidrule{9-9}
      Und && hat && ihm Anna einen Ball && geschenkt && zum Geburtstag? \\
      Oder && schenk && ihm einen Ball! && -- && -- \\ 
    \end{tabular}
  }
  \caption{Feldermodell für Verberstsätze mit Anschlussfeld (AF)}
  \label{fig:VerberstfeldermitAF}
\end{figure}

Mit Verberstsätzen werden neben Entscheidungsfragesätzen und\is{Imperativsatz} Imperativsätzen auch Ausrufe- (\ref{ea:AusrufV1}) und konjunktivische Wunsch- (\ref{ex:WunschV1}) bzw. Bedingungssätze (\ref{ex:BedingungV1}) gebildet.

\ea \label{ea:weitererSatzartenV1}
\ea \label{ea:AusrufV1} Hast du ein Glück!
\ex \label{ex:WunschV1} Hätte ich doch für die Prüfung gelernt!
\ex \label{ex:BedingungV1} Sollte es noch offene Fragen geben, können Sie sich jederzeit bei uns melden.
\z 
\z 

In Dialogsituationen gibt es Sätze, die aussehen wie Verberstsätze, aber dennoch zu den Verbzweitsätzen gehören. Das Beispiel in (\ref{ea:TopicDrop}) ist einer Dialogsequenz\footnote{Tanja Dückers: \textit{Spielzone}. Berlin: Aufbau, 2002 [1999], S. 25.} entnommen. 

\ea \label{ea:TopicDrop} 
\ea \label{ea:HatSiedir} Hat sie dir das nicht erzählt?
\ex \label{ex:Habichwohl} Hab ich wohl gestern nicht mehr mitgekriegt, denke ich.
\z 
\z 

Während die Entscheidungsfrage in (\ref{ea:HatSiedir}) ein Verberstsatz ist, handelt es sich bei der Antwort in (\ref{ex:Habichwohl}) um eine Auslassung der Konstituente im Vorfeld. Ausgelassen wird, was im Gespräch als bekannt bzw. anwesend vorausgesetzt wird. In Abschnitt~\ref{subsubsection:TestNOT} gehen wir auf solche Auslassungen unter dem Stichwort\is{situative Ellipse} \textit{situative Ellipse} ein.


\section{Verbletztsätze}\label{section:Verbletztsatz}

Den dritten Satztyp bilden\is{Verbletztsatz} Verbletztsätze. Wie der Name sagt, steht das Verb, genauer das finite Verb, an letzter Stelle. Verbletztsätze sind abhängige Sätze. Es handelt sich um Nebensätze, die in einen übergeordneten Satz (\textit{Matrixsatz}) eingebettet sind (\ref{ea:Verbletztfunktionen}). In Bezug auf das Verb im einbettenden Satz handelt es sich um Ergänzungen (\ref{ea:VerbletztErgänzung}) oder Angaben (\ref{ex:VerbletztAngabe}). Als Attribute (vgl. Abschnitt~\ref{section:SatzgliederundAttribute}) hängen Nebensätze typischerweise von einem Substantiv ab und sind Teil einer NP (vgl. Abschnitt~\ref{section:NP}) wie in  (\ref{ex:VerbletztAttribut}). 

\ea \label{ea:Verbletztfunktionen}
\ea \label{ea:VerbletztErgänzung} \textit{Dass Anna zu Besuch kommt}, freut uns riesig.
\ex \label{ex:VerbletztAngabe} Wir sind froh, \textit{weil Anna zu Besuch kommt}.
\ex \label{ex:VerbletztAttribut} Die Nachricht, \textit{dass Anna zu Besuch kommt}, hat uns gefreut. 
\z 
\z 

Wir haben bereits gesehen, dass Verbletztsätze als Ergänzungen oder Angaben typischerweise im VF oder im NF von Verbzweit bzw. im NF von Verberstsätzen stehen. Nebensätze in Attributfunktion teilen sich mit dem Substantiv, von dem sie abhängen, ein Feld oder werden ins NF ausgeklammert. 

Intern verfügen auch Verbletztsätze über Felder. Sie zeichnen sich dadurch aus, dass sie durch eine Subjunktion wie \textit{dass} (vgl. Abschnitt~\ref{subsubsection:Subjunktion}) eingeleitet werden und dass das finite Verb ganz am Ende steht, also bei mehrteiligen Verben nach den infiniten Verbteilen (\ref{ea:Verbkomplex_VK}).

\ea \label{ea:Verbkomplex_VK} Anna weiß, \textit{dass die Frist verlängert worden ist}.
\z 

Deshalb sprechen wir erst einmal nicht mehr von der rechten Satzklammer, sondern von einem\is{Verbletztsatz!Verbalkomplex} Verbalkomplex (VK). Zwischen der Subjunktion und dem Verbalkomplex lässt sich auch von einem Mittelfeld sprechen. In diesem gelten die unter \ref{subsection:Mittelfeld} beschriebenen Abfolgetendenzen. Auch Verbletztsätze verfügen über ein Nachfeld. Das Feldermodell für Verbletztsätze wird in der Abbildung~\ref{fig:Verbletztfelder} dargestellt.

\begin{figure}[!htbp]
  \centering
  \oneline{
  \begin{tabular}{cp{0.1em}cp{0.1em}cp{0.1em}c}
    Subjunktion && MF && VK && NF \\
    \cmidrule{1-1}\cmidrule{3-3}\cmidrule{5-5}\cmidrule{7-7}
    dass && Anna uns && besucht && in den Ferien \\
    dass && die Frist && verlängert worden ist && -- \\
  \end{tabular}
  }
  \caption{Feldermodell für Verbletztsätze}
  \label{fig:Verbletztfelder}
\end{figure}

Die Position für die Subjunktion ist obligatorisch zu besetzen, wenn es ein finites Verb gibt. Bei den bereits vorgestellten \textit{zu}-Infinitiven oder infiniten Aufforderungen (sogenannter \textit{Schaffner}imperativ wie \textit{Von der Bahnsteigkante zurücktreten}) kann die Position für die Subjunktion frei bleiben. Bei Angabesätzen mit \textit{um} oder \textit{ohne} und \textit{zu}-Infinitiv stehen \textit{um} und \textit{ohne} auf der Position für Subjunktionen, was in \ref{fig:VerbletztmitohneSubjunktion} dargestellt wird.

\begin{figure}[!htbp]
  \centering
  %\resizebox{\textwidth}{!}{
  \oneline{
    \begin{tabular}{cp{0.1em}cp{0.1em}cp{0.1em}cp{0.1em}c}
       && Subjunktion && MF && VK && NF \\
      \cmidrule{3-3}\cmidrule{5-5}\cmidrule{7-7}\cmidrule{9-9}
      (Sie hat vor,) && -- && an die Ostsee && zu fahren. && -- \\
        && -- && Türen && schließen! && -- \\ 
      (Paul kocht,) && um && Anna eine Freude && zu machen. && -- \\
      (Paul kocht nie,) && ohne && Radio && zu hören. && -- \\
    \end{tabular}
  }
  \caption{Verbletztsätze mit und ohne Subjunktion}
  \label{fig:VerbletztmitohneSubjunktion}
\end{figure}

Verbletztsätze werden auch als selbstständige\is{Verbletztsatz} Ausruf- (\ref{ea:VerbletztAusruf}), konjunktivische Wunsch- (\ref{ex:VerbletztWunsch}), Aufforderungs- (\ref{ex:VerbletztAufforderung}) und Entscheidungsfragesätze (\ref{ex:VerbletztEntscheidungsfrage}) gebraucht. Dies ist durch Auslassungen (Ellipsen) von Prädikaten motiviert, die diese Nebensätze als Ergänzungen regieren wie (\textit{ich weiß nicht,}) \textit{was du wieder hast}. 

\ea \label{ea:VerbletztmitAF}
\ea \label{ea:VerbletztAusruf} Was du wieder hast!
\ex \label{ex:VerbletztWunsch} Wenn ich doch mehr Zeit hätte!
\ex \label{ex:VerbletztAufforderung} Dass du mir ja pünktlich kommst!
\ex \label{ex:VerbletztEntscheidungsfrage} Ob Paul wohl pünktlich kommt?
\z
\z 

Auch Verbletztsätze verfügen über ein Anschlussfeld\is{Anschlussfeld}. Dieses kann durch Konjunktionen besetzt werden, \zB wenn zwei Verbletztsätze wie in (\ref{ea:dassundass}) koordiniert werden. Auch bei selbstständig gebrauchten Verbletztsätzen wie in (\ref{ea:VerbletztmitAF}) kann das Anschlussfeld durch eine Konjunktion besetzt werden wie in (\ref{ex:aberob}).

\ea
\ea \label{ea:dassundass} Die Kandidatin hat gezeigt, dass sie das Amt will \textit{und} dass sie kämpfen kann. 
\ex \label{ex:aberob} \textit{Aber} ob Paul wohl pünktlich kommt?
\z 
\z 

In der Abbildung~\ref{fig:VerbletztfeldermitAF} erweitern wir das Modell des Verbletztsatzes um das Anschlussfeld.

\begin{figure}[!htbp]
  \centering
  %\resizebox{\textwidth}{!}{
  \oneline{
    \begin{tabular}{cp{0.1em}cp{0.1em}cp{0.1em}cp{0.1em}c}
      AF && Subjunktion && MF && VK && NF \\
      \cmidrule{1-1}\cmidrule{3-3}\cmidrule{5-5}\cmidrule{7-7}\cmidrule{9-9}
      und && dass && sie && kämpfen kann && mit aller Kraft \\
      aber && ob && Paul wohl pünktlich && kommt && -- \\ 
    \end{tabular}
  }
  \caption{Feldermodell für Verbletztsätze mit Anschlussfeld (AF)}
  \label{fig:VerbletztfeldermitAF}
\end{figure}

Einen umstrittenen Fall stellen die in\is{Relativsatz} Relativsätzen gebrauchten Pronomen dar, die auf den ersten Blick genauso wie Subjunktionen Verbletztsätze einleiten. Aus diesem Grund verorten sie \citet[280]{Hoehle19862018} und \citet[414]{Eisenberg2020_Satz} auch auf der Position für Subjunktionen. Das Gleiche gilt dann auch für in Relativsätzen und Frage-Nebensätzen gebrauchte Adverbien wie \textit{wohin}, was in \ref{fig:VerbletztRelativSubjunktion} dargestellt ist. 

\begin{figure}[!htbp]
  \centering
  %\resizebox{\textwidth}{!}{
  \oneline{
    \begin{tabular}{cp{0.1em}cp{0.1em}cp{0.1em}cp{0.1em}c}
       && Subjunktion && MF && VK && NF \\
      \cmidrule{3-3}\cmidrule{5-5}\cmidrule{7-7}\cmidrule{9-9}
      %(\textit{Sie sagt},) && \textit{dass} && \textit{sie ans Westufer} && \textit{fährt}. && -- \\
        (Sie mag die Kiefern,) && die && am Westufer && stehen. && -- \\ 
        (Sie weiß nicht,) && wohin && sie sonst && fahren wollte. && -- \\ 
    \end{tabular}
  }
  \caption{Relativsatzeinleiter auf der Subjunktionsposition}
  \label{fig:VerbletztRelativSubjunktion}
\end{figure}

Andere Autoren wie \citet[25--26, 32--33]{Woellstein2014} und \citet[396--397]{Schaefer2018_3Aufl} nehmen an, dass auch Verbletztsätze ein VF haben. Sie verorten die Einleiter des Relativsatzes im Vorfeld des Verbletztsatzes, also nicht auf der Position für Subjunktionen. Hierfür gibt es zwei Argumente. 

Das\is{Relativsatz} erste Argument für die Verortung der Relativsatzeinleiter im Vorfeld liefern Dialektdaten. So führt die \citet[50, §25]{Duden2022} Beispiele aus süddeutschen Dialekten auf, in denen nach dem Relativsatzeinleiter im VF noch die Subjunktion \textit{dass} auf der Position für Subjunktionen steht (\ref{ea:Intmitdass}). Ebenso dialektal ist der Gebrauch des einleitenden Pronomens im VF mit folgendem \textit{wo} in der Position für Subjunktionen (\ref{ex:Relmitwo}).

\ea \label{ea:RelIntimVF}
\ea \label{ea:Intmitdass} Kommt drauf an, mit wem dass sie zu tun haben.
\ex \label{ex:Relmitwo} Du bist der beste Sänger, den wo ich kenne.
\z 
\z 

Zweitens sind die\is{Relativsatz} Relativsatzeinleiter, \textit{w}-Pronomen und entsprechende Frageadverbien (\wzB \textit{wo}, \textit{wann}) Phrasen, denn sie stehen ja für Phrasen. Deutlich zutage tritt der phrasale Charakter bei Fragenebensätzen, wenn das Pronomen wie in (\ref{ex:wer_von_Phrase}) ausgebaut wird.

\ea \label{ea:Fragenebensatz}
\ea \label{ea:wer_Nebensatz} Sie fragt, \textit{wer} ihr helfen könne.
\ex \label{ex:wer_von_Phrase} Sie fragt, \textit{wer von den Anwesenden} ihr helfen könne.
\z 
\z 

Subjunktionen sind hingegen keine Phrasen. Sie regieren den kompletten Nebensatz. Subjunktionen und Einleiter von Relativ- und Fragenebensätzen gehören unterschiedlichen Kategorien an. Auch sind die Funktionen im Satz unterschiedlich: Subjunktionen sind keine Satzglieder. Die im Relativsatz und in Fragenebensätzen gebrauchten Pronomen und Adverbien sind Satzglieder.

Entweder muss angenommen werden, dass auch phrasale Relativsatzeinleiter und Fragepronomen sowie Adverbien wie in der Abbildung~\ref{fig:VerbletztRelativSubjunktion} zu den Subjunktionen zählen, oder aber sie werden im Vorfeld des Verbletztsatzes verortet, dessen Position für die Subjunktion dann leer bleibt. Für beide Analysen gibt es gute Gründe.

Für eine Grammatik in der Schule heißt das: das Feldermodell für Verbletztsätze sollte nicht anhand von Relativsätzen thematisiert werden. Aus sprachpraktischer Perspektive können die Einleiter von Relativ- und indirekten Fragesätzen analog zu Subjunktionen behandelt werden.    



\subsection{Kohärente versus inkohärente Konstruktion}\label{subsection:KohärenteKonstruktion}

In Abschnitt~\ref{section:KonzeptVerbvalenz} haben wir Verben wie \textit{schreiben} und \textit{bitten} kennengelernt, die als Ergänzung satzwertige \textit{zu}-Infinitivkonstruktionen haben. Diese Verben haben wir als Kontrollverben bezeichnet. Denn ihr Subjekt oder ihr Objekt kontrolliert die nicht innerhalb der Infinitivkonstruktion ausgedrückte erste Ergänzung des Infinitivs. Diese Kontrollverben schließen die gesamte \textit{zu}-Infinitivkonstruktion wie in (\ref{ea:bitten_Inkohärent}) inkohärent oder wie in (\ref{ex:bitten_Kohärent}) optional kohärent an. Dies wird in Verbletztsätzen deutlich.

\ea 
\ea \label{ea:bitten_Inkohärent} dass die große Anna den kleinen Paul bittet, [ihr zu helfen]
\ex \label{ex:bitten_Kohärent} dass die große Anna den kleinen Paul, [ihr zu helfen], bittet
\z 
\z

In der inkohärenten Konstruktion (\ref{ea:bitten_Inkohärent}) besetzt die gesamte Infinitivkonstruktion \textit{ihr zu helfen} das Nachfeld, d.\,h. sie steht hinter dem Finitum \textit{bittet}. In der kohärenten Konstruktion (\ref{ex:bitten_Kohärent}) steht dieselbe \textit{zu}-Infinitivkonstruktion im Mittelfeld, d.\,h. vor dem sie regierenden Verb \textit{bittet}.\footnote{Eine Folge der kohärenten Konstruktion ist, dass die Wirkdomäne (\textit{Skopus}) der Negation mehrdeutig ist. Entweder betrifft die Negation das einbettende Verb oder den eingebetteten \textit{zu}-Infinitiv: \textit{dass Anna ihn nicht, ihr zu helfen, bittet}: \gq{Anna bittet ihn nicht, ihr zu helfen} oder \gq{Anna bittet ihn, ihr nicht zu helfen}. Diese Mehrdeutigkeit des Skopus besteht nicht in der inkohärenten Konstruktion: \textit{dass Anna ihn nicht bittet, ihr zu helfen} hat nur eine Lesart. Die andere Lesart verlangt die Negation innerhalb der zu-Infinitivgruppe: \textit{dass Anna ihn bittet, ihr nicht zu helfen.}} In der Abbildung~\ref{fig:In_Kohärente_Konstruktion} wird die Besetzung im Feldermodell gezeigt.

\begin{figure}[!htbp]
  \centering
  %\resizebox{\textwidth}{!}{
  \oneline{
  \begin{tabular}{cp{0.1em}cp{0.1em}cp{0.1em}c}
    Subjunktion && MF && VK && NF \\
    \cmidrule{1-1}\cmidrule{3-3}\cmidrule{5-5}\cmidrule{7-7}
    dass && die große Anna den kleinen Paul && bittet, && ihr zu helfen \\
    dass && die große Anna den kleinen Paul, ihr zu helfen, && bittet && -- \\
  \end{tabular}
  }
  \caption{Inkohärente und kohärente Konstruktion}
  \label{fig:In_Kohärente_Konstruktion}
\end{figure} 

Die komplexe \textit{zu}-Infinitivkonstruktion, hier \textit{ihr zu helfen}, besetzt nicht die Verbalfelder. Im Verbletztsatz steht sie kohärent im Mittelfeld oder inkohärent im Nachfeld. Das Gleiche gilt für Verberst- und Verbzweitsätze. Der \textit{zu}-Infinitiv bildet nicht die rechte Satzklammer.  

Die Vollverben \textit{helfen}, \textit{lernen} und \textit{lehren} erlauben sowohl den Anschluss eines \textit{zu}-Infinitivs (\ref{ea:helfen_zu_Inf}) als auch eines reinen Infinitivs (\ref{ex:helfen_Inf}). 

\ea \label{ea:helfen_(zu)_Inf}
\ea \label{ea:helfen_zu_Inf} Er half ihm zu überleben.\footnote{Der Tagesspiegel, 19.07.2001.}
\ex \label{ex:helfen_Inf} Er hilft uns überleben.\footnote{Berliner Zeitung, 07.05.2005.}
\z 
\z 

Der Anschluss eines Infinitivs ohne \textit{zu} findet sich häufiger, wenn der Infinitiv wie in (\ref{ea:helfen_(zu)_Inf}) nicht durch Angaben oder Ergänzungen erweitert ist. Wenn der Infinitiv kohärent vor dem ihn regierenden Verb steht, wird er wie in (\ref{ea:lernen_lehren_(zu)_Inf}) meist ohne \textit{zu} angeschlossen – unabhängig davon, ob er erweitert (\ref{ea:lernen _Inf}) oder einfach (\ref{ex:lehren_Inf}) ist.


\ea \label{ea:lernen_lehren_(zu)_Inf}
\ea \label{ea:lernen _Inf}  Später, im Universal-Treppenhaus demonstriert er kurz, dass er mal Gesangsunterricht genommen und dabei sogar Arien singen gelernt hat.\footnote{Berliner Zeitung, 15.09.2003.}
\ex \label{ex:lehren_Inf}  Der Fuchs überlässt dem Löwen daraufhin bis auf Weniges seinen Anteil, worauf der Löwe ihn schmunzelnd fragt, wer ihn so schön teilen gelehrt habe.\footnote{Fabel vom Löwenanteil. In: Wikipedia: Die freie Enzyklopädie. 05.10.2024.}
\z 
\z 

Wenn der Infinitiv bzw. die Infinitivgruppe inkohärent im Nachfeld steht, dann wird der \textit{zu}-Infinitiv gebraucht. Das illustrieren wir in (\ref{ea:lernen_lehren_(zu)_Inf2}) anhand der Umstellung der Belege in (\ref{ea:lernen_lehren_(zu)_Inf}). 

\ea \label{ea:lernen_lehren_(zu)_Inf2}
\ea \label{ea:lernen _Inf2}  Später, im Universal-Treppenhaus demonstriert er kurz, dass er mal Gesangsunterricht genommen und dabei sogar gelernt hat, Arien zu singen.
\ex \label{ex:lehren_Inf2}  Der Fuchs überlässt dem Löwen daraufhin bis auf Weniges seinen Anteil, worauf der Löwe ihn schmunzelnd fragt, wer ihn so schön gelehrt habe, zu teilen.
\z 
\z 


Modalverben verhalten sich anders. Wie in Abschnitt~\ref{subsection:Modalverb_INF} erwähnt, sind sie keine Hauptvalenzträger. Sie haben keine komplexen \textit{zu}-Infinitivkonstruktionen als Ergänzung. Modalverben verbinden sich mit einem reinen Infinitiv ohne \textit{zu} zu einem Verbalkomplex. Dieser Verbalkomplex übernimmt die Ergänzungen des infiniten Vollverbs.  Aus der Bildung des Verbalkomplexes aus Modalverb und Infinitiv folgt, dass nur eine kohärente Konstruktionsweise wie in (\ref{ea:sollen_kohärent}) möglich ist. Eine inkohärente Konstruktion wie in (\ref{ex:sollen_Inkohärent}) ist nicht möglich.

\ea 
\ea [] {dass die große Anna dem kleinen Paul helfen soll}\label{ea:sollen_kohärent}
\ex [*] {dass die große Anna soll dem kleinen Paul helfen}\label{ex:sollen_Inkohärent}
\z 
\z
 
Das Modalverb \textit{soll} verbindet sich mit dem zweiwertigen \textit{helfen} zu dem komplexen Prädikat \textit{helfen soll}. Der Komplex verfügt über zwei Ergänzungen, nämlich \textit{die große Anna} und \textit{dem kleinen Paul}. Der pure Infinitiv \textit{helfen} verbindet sich  nicht mit seiner ursprünglich zweiten Ergänzung zu der komplexen Infinitivkonstruktion \textit{dem kleinen Paul helfen}. Er ist nicht satzwertig. Folglich kann \textit{dem kleinen Paul helfen} auch nicht inkohärent im Nachfeld stehen. In Abbildung~\ref{fig:Kohärente_MV} zeigen wir die Besetzung im Feldermodell.

\begin{figure}[!htbp]
  \centering
  %\resizebox{\textwidth}{!}{
  \oneline{
  \begin{tabular}{cp{0.1em}cp{0.1em}cp{0.1em}c}
    Subjunktion && MF && VK && NF \\
    \cmidrule{1-1}\cmidrule{3-3}\cmidrule{5-5}\cmidrule{7-7}
    dass && die große Anna dem kleinen Paul && helfen soll && -- \\
  \end{tabular}
  }
  \caption{Obligatorisch kohärent konstruierende Modalverben}
  \label{fig:Kohärente_MV}
\end{figure} 

Der Infinitiv besetzt zusammen mit dem Modalverb im Verbletztsatz die Stelle des Verbalkomplexes (VK). In Verberst- und Verbzweitsätzen bildet das finite Modalverb die linke Satzklammer und der Infinitiv die rechte Satzklammer.

Wie die Modalverben konstruieren auch alle Hilfsverben obligatorisch kohärent. So verbindet sich zuerst ein Hilfsverb wie \textit{hat} in (\ref{ea:hat_geholfen_kohärent}) mit dem Partizip II \textit{geholfen} zu dem Verbalkomplex \textit{geholfen hat}. Dieser Verbalkomplex übernimmt die Ergänzungen des Vollverbs \textit{helfen}. Das Partizip II kann nicht alleine mit der zweiten Ergänzung \textit{dem kleinen Paul} im Nachfeld stehen. Die kohärente Konstruktion wie in (\ref{ea:hat_geholfen_kohärent}) ist obligatorisch, die inkohärente wie in (\ref{ex:hat_geholfen_inkohärent}) hingegen ungrammatisch. 

\ea 
\ea [] {dass die große Anna dem kleinen Paul geholfen hat}\label{ea:hat_geholfen_kohärent}
\ex [*] {dass die große Anna hat dem kleinen Paul geholfen}\label{ex:hat_geholfen_inkohärent}
\z 
\z
 
Das Hilfsverb bildet zusammen mit dem Partizip II ein komplexes Prädikat. Dieses besetzt im Verbletztsatz die Stelle des Verbalkomplexes (VK). In Verberst- und Verbletztsätzen bildet das finite Hilfsverb die linke Satzklammer und das Partizip II die rechte Satzklammer.

Obligatorisch kohärent konstruieren auch AcI-Verben wie \textit{sehen} (vgl. Abschnitt~\ref{subsection:AcI_Inf}) in (\ref{ea:tanzen_sieht_kohärent}). Auch sie bilden zusammen mit einem Infinitiv wie \textit{tanzen} ein komplexes Prädikat: \textit{tanzen sieht}. Der Infinitiv \textit{tanzen} alleine bildet keine Einheit mit der zweiten Ergänzung \textit{den kleinen Paul}. Es bleibt nur die kohärente Konstruktion wie in (\ref{ea:tanzen_sieht_kohärent}). Die inkohärente Konstruktion, d.\,h. die Auslagerung von \textit{den kleinen Paul tanzen} ins Nachfeld wie in (\ref{ex:tanzen_sieht_inkohärent}) ist nicht möglich.

\ea 
\ea  [] {dass die große Anna den kleinen Paul tanzen sieht}\label{ea:tanzen_sieht_kohärent}
\ex [*] {dass die große Anna sieht den kleinen Paul tanzen}\label{ex:tanzen_sieht_inkohärent}
\z 
\z

Das AcI-Verb bildet zusammen mit dem Infinitiv ein komplexes Prädikat. Dieses besetzt im Verbletztsatz die Stelle des Verbalkomplexes (VK). In Verberst- und Verbletztsätzen bildet das finite AcI-Verb die linke Satzklammer und der Infinitiv die rechte Satzklammer.

In aller Regel konstruieren auch Modalitätsverben (vgl. Abschnitt~\ref{subsection:Modalitätsverb_zu_Inf}) kohärent. Im Unterschied zu den Modalverben verlangen die Modalitätsverben wie \textit{scheinen} einen \textit{zu}-Infinitiv wie \textit{zu helfen}. Zusammen bilden sie ein komplexes Prädikat wie \textit{zu helfen scheint}. Das komplexe Prädikat übernimmt die Ergänzungen des Vollverbs. Die Konstruktion ist kohärent (\ref{ea:zu_helfen_scheint_kohärent}). Der \textit{zu}-Infinitiv kann nicht mit seiner ursprünglichen Ergänzung \textit{dem kleinen Paul} im Nachfeld stehen (\ref{ex:zu_helfen_scheint_inkohärent}). 

\ea 
\ea  [] {dass die große Anna dem kleinen Paul zu helfen scheint}\label{ea:zu_helfen_scheint_kohärent}
\ex [*] {dass die große Anna scheint dem kleinen Paul zu helfen}\label{ex:zu_helfen_scheint_inkohärent}
\z 
\z

Das Modalitätsverb bildet zusammen mit dem \textit{zu}-Infinitiv ein komplexes Prädikat. Dieses besetzt im Verbletztsatz die Stelle des Verbalkomplexes (VK). In Verberst- und Verbletztsätzen bildet das finite Modalitätsverb die linke Satzklammer und der \textit{zu}-Infinitiv die rechte Satzklammer.\footnote{\citet[1283]{ZifHoffStr1997} zeigen, dass \textit{drohen} inkohärent konstruieren kann: \textit{weil die Situation noch immer drohte, recht unübersichtlich zu werden}. Auch mit \textit{versprechen} finden sich inkohärent konstruierte Belege: \textit{die große China-Ausstellung 1998, die verspricht, ein besonderes Ereignis zu werden.} Der Beleg stammt aus \citet[137]{Reis2005}, die diese Herausstellung als einen besonderen \textit{3. Konstruktionstyp} trotz kohärenter Konstruktion bezeichnet. Diese besondere Konstruktionsform spricht für eine weniger starke Grammatikalisierung, die sich in doppeldeutigen Lesarten zeigt: \textit{Paul droht}(,) \textit{seinen Komplizen zu verraten}/\textit{Paul verspricht}(,) \textit{ein großer Schauspieler zu werden.}}

Für die obligatorisch kohärent konstruierenden Hilfsverben, AcI-Verben und Modalitätsverben wird in der Abbildung~\ref{fig:Kohärente_Konstruktionen} die Felderbesetzung dargestellt.

\begin{figure}[!htbp]
  \centering
  %\resizebox{\textwidth}{!}{
  \oneline{
  \begin{tabular}{cp{0.1em}cp{0.1em}cp{0.1em}c}
    Subjunktion && MF && VK && NF \\
    \cmidrule{1-1}\cmidrule{3-3}\cmidrule{5-5}\cmidrule{7-7}
    dass && die große Anna dem kleinen Paul && geholfen hat && -- \\
    dass && die große Anna den kleinen Paul && tanzen sieht && -- \\
    dass && die große Anna dem kleinen Paul && zu helfen scheint && -- \\
  \end{tabular}
  }
  \caption{Obligatorisch kohärent konstruierende Verben}
  \label{fig:Kohärente_Konstruktionen}
\end{figure} 

\newpage
\subsection{Verbalkomplex}\label{subsection:RechteSatzklammer}

Bei Verbletztsätzen\is{Verbletztsatz!Verbalkomplex} steht der gesamte Verbalkomplex am Ende. Der wesentliche Unterschied zu Verbzweit- und Verberstsätzen besteht darin, dass das Finitum hinter den infiniten Verbteilen steht. Wenn das Verb einteilig ist, so steht das Finitum alleine in der Position für den Verbalkomplex.

In Abschnitt~\ref{section:RegierteVerben} sind wir unter dem Begriff der Statusrektion schon auf die Stellung mehrerer Verbteile im Verbalkomplex eingegangen. Wir haben gezeigt, dass die Abfolge der Verbteile umgekehrt zu den Rektionsverhältnissen verläuft, also vom Ende zum Anfang des Verbalkomplexes (\ref{ea:Unterfeld}). 

 \ea \label{ea:Unterfeld}
 \ea \label{ea:2InfinitiveRSK} dass Anna den Jungen singen $\leftarrow$ gehört $\leftarrow$ hat
 \ex \label{ex:2InfinitiveRSK2} dass Anna den Jungen singen $\leftarrow$ hören $\leftarrow$ können $\leftarrow$ wird
 \z 
 \z 

In (\ref{ea:2InfinitiveRSK}) fordert das Perfekt-Hilfsverb \textit{hat} ein Partizip II, hier \textit{gehört}. Das als AcI-Verb gebrauchte \textit{gehört} verlangt seinerseits nach einem Infinitiv, hier \textit{singen}.

In (\ref{ex:2InfinitiveRSK2}) verlangt das Futur-Hilfsverb \textit{wird} nach einem Infinitiv, hier \textit{können}. Das Modalverb \textit{können} verlangt nach einem Infinitiv, hier \textit{hören}. Das als AcI-Verb gebrauchte \textit{hören} verlangt nach einem Infinitiv, hier \textit{singen}.
 
 Diese Abfolge gilt jedoch nicht immer. Unter bestimmten Umständen steht das Finitum und wie in (\ref{ex:OF}) eine von diesem regierte Verbform nicht am Ende des Verbalkomplexes, sondern an dessen Beginn wie in (\ref{ea:EIinRSK}) und (\ref{ex:OF}). 

 \ea \label{ea:OFUF}
 \ea \label{ea:EIinRSK} dass Anna ihn hat singen hören
 \ex \label{ex:OF} dass Anna ihn wird können singen hören
 \z 
 \z 
 
 Zur Beschreibung dieser Abfolge unterteilt man den Verbalkomplex in\is{Verbletztsatz!Oberfeld} \textit{Oberfeld} (OF) und\is{Verbletztsatz!Unterfeld} \textit{Unterfeld} (UF), was in Abbildung~\ref{fig:OberfeldUnterfeldRSK} dargestellt wird. Ist das OF komplex besetzt, so folgt die regierte Verbform auf die regierende, also umgekehrt zur Stellung im UF.
 
 \begin{figure}[!htbp]
  \centering
  %\resizebox{\textwidth}{!}{
  \oneline{
  \begin{tabular}{cp{0.1em}cp{0.1em}cp{0.1em}c}
    Subjunktion && MF && VK && \\
    \cmidrule{1-1}\cmidrule{3-3}\cmidrule{5-5}
    dass && Anna ihn && []\textsubscript{\textsc{of}} [singen gehört hat]\textsubscript{\textsc{uf}} \\
    dass && Anna ihn && [hat]\textsubscript{\textsc{of}} [singen hören]\textsubscript{\textsc{uf}}\\
    dass && Anna ihn && [wird können]\textsubscript{\textsc{of}} [singen hören]\textsubscript{\textsc{uf}} 
  \end{tabular}
  }
  \caption{Ober- und Unterfeld in der RSK}
  \label{fig:OberfeldUnterfeldRSK}
\end{figure}

Das Finitum kann innerhalb des VK ins Oberfeld vorangestellt werden, wenn es wie in (\ref{ea:2InfinitiveRSK}) mindestens zwei infinite Verben gibt. Mit der Voranstellung des Finitums ins\is{Verbletztsatz!Oberfeld} Oberfeld geht einher, dass nicht das Partizip II des regierten Verbs (\textit{gehört}), sondern ein sogenannter\is{Ersatzinfinitiv} \textit{Ersatzinfinitiv} (Ersatz für das Partizip II) gebraucht wird. Statt \textit{singen gehört hat} heißt es dann \textit{hat singen hören}. 

Wenn das Finitum ein Hilfsverb für die Perfektbildung ist, das ein Modalverb regiert (\textit{gewollt hat}), welches seinerseits einen Infinitiv regiert, liegt ein Sonderfall vor. In diesem Fall gelten die\is{Verbletztsatz!Oberfeld} Oberfeldbesetzung des Finitums (\textit{hat singen wollen} statt \textit{wollen singen hat}) und der Gebrauch des\is{Ersatzinfinitiv} Ersatzinfinitivs (\textit{wollen} statt \textit{gewollt}) als obligatorisch (\ref{ea:OberfeldobligatorischAuxINFMod}).\footnote{Nach \citet[64]{Bech1983} kommen Verbalkomplexe mit mehr als vier Verben nur sehr selten vor. Er konstruiert ein Beispiel mit fünf Verben: \textit{ob man ihn hier} [\textit{wird können}]\textsubscript{OF} [\textit{liegen bleiben lassen}]\textsubscript{UF}. Bei weiterem Interesse verweisen wir auf \citet[67--68]{Pafel2011}.}

\ea \label{ea:Oberfeldobligatorisch} 
\ea [] {dass Anna den Jungen hat singen hören wollen}\label{ea:OberfeldobligatorischAuxINFMod}
\ex [*] {dass Anna den Jungen singen hören wollen/gewollt hat}\label{ex:Oberfeldverstoß}
\z 
\z 

Auch in Verbzweit- und Verberstsätzen wird bei den AcI-Verben \textit{sehen} und \textit{hören} häufig optional der\is{Ersatzinfinitiv} Ersatzinfinitiv im Perfekt gebraucht (\ref{ea:EIbeiAcI}). Beim kausativ gebrauchten \textit{lassen} (\ref{ex:lassenEI}) und bei den Modalverben ist der Ersatzinfinitiv im Perfekt obligatorisch (\ref{ex:EIMV}). Das Partizip II muss gebraucht werden, wenn die Modalverben als Vollverben mit einer NP im Akkusativ gebraucht werden (\ref{ex:PIIMV}).

\ea \label{ea:Ersatzinfinitiv}
\ea \label{ea:EIbeiAcI} Sie hat ihn singen hören/gehört.
\ex \label{ex:lassenEI} Er hat die Kinder das Zimmer aufräumen lassen/*gelassen.
\ex \label{ex:EIMV} Hat er an die Ostsee fahren wollen/*gewollt?
\ex \label{ex:PIIMV} Er hat ein Eis gewollt/*wollen.
\z 
\z 

Warum insbesondere bei den Modalverben obligatorisch der\is{Ersatzinfinitiv} Ersatzinfinitiv gebraucht wird und dies mit der Oberfeldbesetzung des Finitums einhergeht, ist weitgehend ungeklärt. Einen interessanten Ansatz stellen \citet{EisenbergSmithTeuber2001} vor. Die Bildung des Ersatzinfinitivs erklären \citet[256]{EisenbergSmithTeuber2001} damit, dass Modalverben keine Nachzustände implizieren. Weder die erste Ergänzung, \zB jemand, der singen will, noch die zweite Ergänzung, \zB ein Lied, das jemand singen will, verändern sich innerhalb eines Singen-wollen-Szenarios bzw. weisen innerhalb von zwei Zeitintervallen verschiedene Merkmale auf. Jene Veränderung im zweiten Zeitabschnitt (Nachzustand) ist aber gerade die Ausdrucksdomäne vom Partizip II. So bezieht sich das Partizip \textit{geöffnet} beim Perfekt auf den Nachzustand der zweiten Ergänzung, deren Referent zuvor geschlossen war, und \textit{eingeschlafen} auf die erste Ergänzung, deren Referent vorher wach war. Da Modalverben keine Nachzustände implizieren, sei demnach die Bildung des Partizip II semantisch blockiert. An seiner Stelle werde der Infinitiv gebraucht.\footnote{Bei Verben wie \textit{schlafen}, die ebenfalls keine Zustandsveränderung implizieren, habe sich, so \citet[256]{EisenbergSmithTeuber2001}, der Druck der Grammatikalisierung des Perfekts analog zu \textit{öffnen} (etwas ist im ersten Zeitintervall offen und im zweiten geschlossen) und \textit{geöffnet haben} durchgesetzt. Modalverben seien auf Grund ihrer flexionsmorphologischen Besonderheit (aus Präteritumformen entstanden zu sein) diesem Druck nicht ausgesetzt gewesen. Ein Indiz hierfür ist, dass es in der Sprachgeschichte früher mehr Verben gab, die mit dem Ersatzinfinitiv gebraucht worden sind, \zB \textit{anfangen}, \textit{pflegen} und (\textit{ver})\textit{suchen}. Bei diesen Verben war aber der Anpassungsdruck so stark, dass sie im heutigen Standard keinen Ersatzinfinitv mehr bilden. In Tirol und in Bayern werden auch heute Ersatzinfinitive analog zu den Modalverben gebraucht, \zB \textit{ich habs mir nicht sagen trauen}. Einen Überblick zur historischen Entwicklung geben \citet[179--193]{Fleischer2001}.} 

Die Oberfeldstellung des Finitums wird auf den\is{Ersatzinfinitiv} Ersatzinfinitiv zurückgeführt. Im Verbletztsatz müsste es analog zu \textit{geschlafen hat} bei Modalverben \zB \textit{wollen hat} heißen. Während aber von \textit{geschlafen} in der Sprachverarbeitung auf das Auxiliar \textit{haben} geschlossen werden kann, ist dies beim Ersatzinfinitiv \textit{wollen} gerade nicht möglich. Es könnte auch heißen \textit{schlafen wollen wird}. Die Voranstellung des Perfektauxiliars im OF, so \citet[258]{EisenbergSmithTeuber2001}, erleichtere die Verarbeitung des Verbalkomplexes.


\section{Differenzierte Feldermodelle und uniformes Grundmodell}\label{section:differenziert_uniform_Feldermodell}

Bisher haben\is{Feldermodell!differenziert} wir in Anlehnung an \citet{Hoehle19862018}, \citet[409, 414]{Eisenberg2020_Satz} und \citet[68--79]{Agel2017} drei verschiedene Feldermodelle vorgestellt, nämlich jeweils eins für Verbzweit-, Verberst- und Verbletztsätze. Diese drei Modelle werden in der Tabelle~\ref{tab:3Feldermodelle} nochmals dargestellt.

\begin{table}
  \centering
  %\resizebox{1\textwidth}{!}{
  \oneline{
    \begin{tabular}{lp{0.3cm}lllllll}
    \lsptoprule
    Satztyp &&  &  & & & & \\
    \midrule
    Verbzweit && AF & VF & LSK\textsubscript{\textsc{finit}} & MF & RSK\textsubscript{\textsc{infinit}} & NF \\
    Verberst && AF & & LSK\textsubscript{\textsc{finit}} & MF & RSK\textsubscript{\textsc{infinit}} & NF \\
    Verbletzt && AF & (VF) & Subjunktion & MF & VK\textsubscript{\textsc{infinit finit}} & NF \\ 
    \lspbottomrule
  \end{tabular}
  }
  \caption{Differenzmodell: 3 Feldermodelle für 3 Satztypen}
  \label{tab:3Feldermodelle}
\end{table}

Die\is{Feldermodell!differenziert} drei verschiedenen Modelle sind fundamental, um die Besonderheiten der verschiedenen Satztypen zu beschreiben. Wer die Felder und ihre Besetzungsregeln nicht intuitiv beherrscht, für den stellen die drei Modelle eine unverzichtbare Grundlage für die Sprech- und Schreibpraxis dar. Deshalb sind für zukünftige Lehrerinnen und Lehrer gerade die Unterschiede zwischen den Satztypen wesentlich. Wer sich die Unterschiede klar gemacht hat, der wird allerdings auch Gemeinsamkeiten entdecken. 

Trotz der\is{Feldermodell!uniform} wesentlichen Unterschiede zwischen den Satztypen fällt in der Tabelle~\ref{tab:3Feldermodelle} auf, dass der Verbalkomplex in Verbletztsätzen an der gleichen Position wie die RSK in Verbzweit- und Verberstsätzen steht. Die infiniten Verbteile in der RSK weisen die gleichen Stellungen wie im Verbalkomplex auf, weshalb beide zu einem Feld zusammengefasst werden können. 

Die LSK-Position für das Finitum entspricht der Position für die Subjunktion in Verbletztsätzen. Diesen Positionsbezug erfasst \citet[26, 33]{Woellstein2014} im \textit{uniformen} Feldermodell. Auch \citet[289]{Hoehle19862018}, der für das in \ref{tab:3Feldermodelle} dargestellte Differenzmodell argumentiert, sieht jenen Positionsbezug und hält einen Satzklammerbegriff aus Finitum oder Subjunktion sowie dem Verbalkomplex in der RSK für möglich.\footnote{\citet[289]{Hoehle19862018} schreibt allerdings: \gqq{Das ist zwar überflüssig, aber soweit nicht schädlich.}} Jener Bezug hat den Vorteil, paralleles Stellungsverhalten zu erfassen. Die LSK ist dann nicht mehr die Stelle für das Finitum, sondern die Position, auf welcher entweder das Finitum oder im Verbletztsatz eine Subjunktion steht. 

Im Rahmen des\is{Feldermodell!uniform} uniformen Modells muss angenommen werden, dass alle Satztypen über ein VF verfügen. In Verberstsätzen bleibt dieses obligatorisch leer. In Verletztsätzen kann es, je nach Analyse, bei Relativ- und Fragenebensätzen besetzt werden. Wir stellen das uniforme Modell in der Tabelle~\ref{tab:uniformesFeldermodell} dar. Da es ein nicht obligatorisch zu besetzendes AF und NF in allen Satztypen gibt, vernachlässigen wir diese Felder der Übersichtlichkeit halber.  

\begin{table}
  \centering
  %\resizebox{1\textwidth}{!}{
  \oneline{
    \begin{tabular}{lp{0.3cm}lllll}
    \lsptoprule
    Satztyp && VF  & LSK & MF & RSK \\
    \midrule
    Verbzweit && 1 Satzkonstituente & Finitum & versch. Konstituenten & infiniter Verbrest &  \\
    Verberst && -- & Finitum & versch. Konstituenten & infiniter Verbrest \\
    Verbletzt &&  (--) & Subjunktion & versch. Konstituenten & gesamter Verbalkomplex \\ 
    \lspbottomrule
  \end{tabular}
  }
  \caption{Uniformes Feldermodell für die 3 Satztypen}
  \label{tab:uniformesFeldermodell}
\end{table}

In Analogie zu Verbzweit- und Verberstsätzen kann auch bei Verbletztsätzen von einer \textit{Klammer} (zwischen Subjunktion und Finitum) gesprochen werden. Hierin liegt aber auch ein Nachteil. Denn Finitum und infiniter Verbrest bilden als eine syntaktische Konstituente durch Distanzstellung tatsächlich eine Klammer. In Verbletztsätzen stehen infiniter Verbrest und Finitum zusammen in der RSK und bilden gerade keine Klammer. Die Subjunktion steht dort, wo in Verbzweit- und Verberstsätzen das Finitum steht. 

Diese Positionsidentität kann der Grund dafür sein, dass sowohl das Finitum als auch die Subjunktion durch Betonung den Wahrheitsgehalt eines Satzes fokussieren. Hierauf weist \citet[148]{Reis1987} hin, von der auch die Sätze in (\ref{ea:LSKBetonungWahrheitsgehalt}) übernommen sind.\footnote{\citet[383]{Hoehle1992_2018} bezeichnet in einer 1992 erstmals veröffentlichten Arbeit den Effekt, welchen die Betonung des Verbs auf den Wahrheitsgehalt hat, als \textit{Verum-Fokus}.} Unabhängig vom Satztyp unterstreicht die Betonung des Elements in der LSK den Wahrheitsgehalt des Satzes. Die Betonung wird durch Großschreibung angezeigt.

\ea \label{ea:LSKBetonungWahrheitsgehalt}
\ea \label{ea:V2LSKBetonung} Er HAT sie umgebracht.
\ex \label{ex:V1LSKBEtonung} HAT er sie nun umgebracht oder nicht?
\ex \label{ex:VletztLSKBetonung} Es kommt heute nicht mehr darauf an, DASS er sie umgebracht hat, sondern wie er dazu heute steht.
\z 
\z 

Die Positionsidentität von Finitum und Subjunktion zeigt sich auch darin, dass \zB im Bairischen Subjunktionen Flexionsmerkmale tragen, welche im Standarddeutschen die Finitheit eines Verbs markieren. Das Beispiel (\ref{ea:SubjflektiertBairisch}) stammt von \citet[207]{Grewendorf1988}. 

\ea \label{ea:SubjflektiertBairisch} I woaß, daß\textit{st}/ob\textit{st} du a Spitzbua bi\textit{st}.
\z


\section{Deutsch als Verbzweitsprache mit Verbletztstellung und freier Wortfolge}\label{section:DeutschV2VletztfreieWortstellung}

Für das Deutsche sind die Zweitstellung des Finitums \textit{und} die Verbletztstellung charakteristisch. Nicht minder wesentlich ist, dass Deutsch als eine Sprache mit relativ freier Wortfolge gilt. So ist der Titel eines Beitrags von \citet{Haftka1996} besonders eindrücklich: \gqq{Deutsch ist eine V/2-Sprache mit Verbendstellung und freier Wortfolge}. Diese typologischen Besonderheiten werden wir abschließend im Kontrast zum Englischen illustrieren.

Dass Deutsch eine\is{Verbzweitsprache} Verbzweitsprache ist, zeigt sich in Aussage- (\ref{ea:V2deutsch}) und Ergänzungsfragesätzen. Aus kommunikativem Blickwinkel handelt es sich um die Grundstellung des deutschen Satzes, worauf auch die Bezeichnung \textit{Kern}satz abzielt. Die Verbzweitstellung tritt deutlich beim Vergleich mit dem Englischen (welches unter den westgermanischen Sprachen eine Ausnahme darstellt) hervor (\ref{ex:nichtV2englisch}).

\ea \label{ea:V2deutschenglisch}
\ea \label{ea:V2deutsch} Gestern \textit{traf} die große Frau den kleinen Mann.
\ex \label{ex:nichtV2englisch} Yesterday, the big woman \textit{met} the little man.
\z 
\z 

Wenn, \zB aus Fokusgründen, keine Konstituente im Vorfeld steht, dann sorgt im Deutschen der\is{es@\textit{es}!Platzhalter} Platzhalter (Expletivum) \textit{es} für die Verbzweitstellung und kennzeichnet somit den Satztyp (\ref{ea:Expletivum}). 

\ea \label{ea:Expletivum} \textit{Es} traf gestern die große Frau den kleinen Mann.
\z 


Die Eigenschaft des Deutschen als\is{Verbletztsprache} Verbletztsprache zeigt sich am augenfälligsten in abhängigen Sätzen (\ref{ea:Vletztdeutsch}). Auch diese Eigenschaft lässt sich kontrastiv zum Englischen darstellen (\ref{ex:nichtVletztenglisch}).

\ea \label{ea:Vletztdeutschenglisch}
\ea \label{ea:Vletztdeutsch} Ich weiß, dass die große Frau gestern den kleinen Mann \textit{traf}.
\ex \label{ex:nichtVletztenglisch} I know, that the big woman \textit{met} the little man yesterday.
\z 
\z 

Diese Eigenschaft zeigt sich aber auch bei Zweit- oder Erststellung des Finitiums, wenn es einen infiniten Verbteil gibt. Das infinite Vollverb steht am Ende (\ref{ea:O_V}), was sich grundlegend vom Englischen unterscheidet (\ref{ex:V_O}). 

\ea \label{ea:VletztbeiFin2}
\ea \label{ea:O_V} Der kleine Mann will die große Frau \textit{treffen}.
\ex \label{ex:V_O} The little man wants \textit{to meet} the big woman.
\z 
\z 

Wir müssen zwei Eigenschaften des deutschen Satzbaus voneinander trennen. Das ist zum einen die dargestellte Zweitstellung des Finitums und zum anderen die Eigenschaft als Verbletztsprache. Unter Verbletztsprache versteht man die grundlegende End- bzw. Letztstellung des regierenden Vollverbs. In diesem Sinne ist Deutsch eine\is{Objekt-Verbsprache} Objekt-Verb-Sprache (OV-Sprache) und Englisch eine Verb-Objekt-Sprache (VO-Sprache). Dieser Unterschied wird bereits beim Vergleich von Infinitivgruppen deutlich. So heißt es im Deutschen \textit{einen Baum pflanzen}, im Englischen hingegen \textit{to plant a tree}.

Auch die Stellung der sogenannten trennbaren Verbpartikeln gilt als Indiz für die\is{Objekt-Verbsprache} strukturelle Verbletztstellung. So stehen die Verbpartikeln nur bei Verbletztstellung zusammen mit dem entsprechenden Verb (\ref{ea:VptVzusammen}). Bei Verbzweit- (\ref{ex:VptV2getrennt}) und Verberststellung stehen Verb und Partikel getrennt und bilden die Satzklammer.

\ea \label{ea:VptVStellung}
\ea \label{ea:VptVzusammen} (Man weiß,) dass Anna ein Theaterstück \textit{aufführt}.
\ex \label{ex:VptV2getrennt} Anna \textit{führt} ein Theaterstück \textit{auf}.
\ex \label{ex:VptV1getrennt} \textit{Führt} Anna ein Theaterstück \textit{auf}?
\z 
\z 

Ein anderes Indiz für die Annahme der\is{Objekt-Verbsprache} Verbletzt- als Grundstellung führt \citet[370--371]{Hoehle19912018} an. Manche Verben werden wie \textit{uraufführen} von einem Substantiv (\textit{Uraufführung}) gebildet.\footnote{\citet[370]{Hoehle19912018} zählt auch \textit{zwischenfinanzieren}, \textit{wettrudern}, \textit{rückfragen} als sogenannte verbale Pseudokomposita auf. Um diese Art der Wortbildung geht es in Abschnitt~\ref{subsection:Rückbildung}.} Sie werden aber fast nur in Verbletztstellung als der originären Stellung gebraucht (\ref{ea:uraufführenRSK}). Die Stellungen in Verbzweit- (\ref{ex:V2führenurauf}) und Verberstsätzen (\ref{ex:V1führenurauf}) sind äußerst selten. 

\ea \label{ea:uraufführen}
\ea [] {(Man weiß,) dass Anna das Theaterstück uraufführt.}\label{ea:uraufführenRSK}
\ex [] {Anna führt das Theaterstück urauf.}\label{ex:V2führenurauf} 
\ex [] {Führt Anna das Theaterstück nicht urauf?}\label{ex:V1führenurauf}
\z 
\z 

Ein weiteres Indiz für die\is{Objekt-Verbsprache} Verbletzt- als Grundstellung des Finitums ist der Wirkungsbereich (\textit{Skopus}) von Ausdrücken im Satz. Der Skopus von Angaben hängt von ihrer Reihenfolge ab. Die links stehende Angabe hat normalerweise Skopus über die folgende Angabe und das Verb in Verbletztstellung (\ref{ea:SkopusVletzt}). Wir übernehmen ein leicht abgewandeltes Beispiel aus \citet[77]{Mueller2010}. 

\ea \label{ea:SkopusVletzt}
\ea \label{ea:absichtlichniemalsVletzt} dass Paul [absichtlich [niemals lacht]]
\ex \label{ex:niemalsabsichltichVletzt} dass Paul [niemals [absichtlich lacht]]
\z 
\z 

Während es in (\ref{ea:absichtlichniemalsVletzt}) darum geht, mit Absicht, niemals zu lachen, wird in (\ref{ex:niemalsabsichltichVletzt}) ausgedrückt, dass niemals mit Vorsatz gelacht wird. Diese Skopusverhältnisse bleiben interessanterweise unverändert, wenn das Verb an zweiter (\ref{ea:absichtlichniemalsV2}) oder erster Stelle (\ref{ex:niemalsabsichltichV1}) steht. Mit anderen Worten, es sieht so aus, als ob die Skopusverhältnisse bei Verbletztstellung als Grundstellung festgelegt werden.

\ea \label{ea:SkopusV12}
\ea \label{ea:absichtlichniemalsV2} Paul lacht absichtlich niemals.
\ex \label{ex:niemalsabsichltichV1} Lacht Paul niemals absichtlich? 
\z 
\z 

Wenn man von der Verbletztstellung als Grundstellung ausgeht, dann lässt sich in einem Schritt aus dieser die Verberststellung ableiten. Betrachten wir hierzu noch einmal einen Verbletztsatz. In (\ref{ea:V_Letzt_Grund}) steht die Subjunktion \textit{dass} in der LSK und eröffnet das MF. Im Mittelfeld stehen die Ergänzungen und Angaben des Verbs: \textit{die große Frau}, \textit{den kleinen Mann} und \textit{gestern}. In der RSK steht der Verbalkomplex \textit{getroffen hat}.

\ea\label{ea:V_Letzt_Grund} dass die große Frau gestern den kleinen Mann getroffen hat
\z 

Wenn es keine Subjunktion in der LSK gibt, dann steht das Finitum \textit{hat} nicht mehr in der RSK, sondern ganz vorne in der leeren LSK. Dass das Finitum statt in der RSK nun in der LSK steht, illustriert der Pfeil in (\ref{ea:V_Erst_1}).    
    
\begingroup
    \tikzstyle{every picture}+=[remember picture, inner sep=0pt, baseline, anchor=base]%
\ea\label{ea:V_Erst_1} {} \tikz\node(wh2){Hat}; die große Frau gestern den kleinen Mann getroffen \tikz\node(wh1){[\phantom{hat}]};	\begin{tikzpicture}[overlay]
	\draw[-latex,rounded corners=.25em](wh1.south)--+(0,-1em)-|(wh2.south)node [anchor=center,pos=0.25,fill=white] {\small\textit{Finitum in LSK statt in RSK}};
	\end{tikzpicture}
 \z 

\vspace{1.25em}

Aus dem Verberstsatz (\ref{ea:V_Erst_1}) können wir in einem weiteren Schritt einen Verbzweitsatz ableiten. Unter Ableitung verstehen wir keine Verschiebungen oder Transformationen, sondern bestimmte Stellungen des Finitums oder einer Satzkonstituente statt einer anderen Stellung. Im Verbzweitsatz gibt es vor der LSK das VF. Das VF muss im Verbzweitsatz von einer Satzkonstituente besetzt sein, \zB von einer Ergänzung oder von einer Angabe. Diese steht im Verbzweitsatz statt im Mittelfeld im Vorfeld. In (\ref{ea:V_Zweit_2}) illustrieren wir das für die Ergänzung \textit{die große Frau} und in (\ref{ea:V_Zweit_3}) für die Angabe \textit{gestern}. 

    \tikzstyle{every picture}+=[remember picture, inner sep=0pt, baseline, anchor=base]%
\ea\label{ea:V_Zweit_2} {} \tikz\node(wh2){Die große Frau}; hat \tikz\node(wh1){[\phantom{die große Frau}]}; gestern den kleinen Mann getroffen	\begin{tikzpicture}[overlay]
	\draw[-latex,rounded corners=.25em](wh1.south)--+(0,-1em)-|(wh2.south)node [anchor=center,pos=0.25,fill=white] {\small\textit{im VF statt im MF}};
	\end{tikzpicture}
\z 

\vspace{1.25em}

    \tikzstyle{every picture}+=[remember picture, inner sep=0pt, baseline, anchor=base]%
\ea\label{ea:V_Zweit_3} {} \tikz\node(wh2){Gestern}; hat die große Frau \tikz\node(wh1){[\phantom{gestern}]}; den kleinen Mann getroffen \begin{tikzpicture}[overlay]
	\draw[-latex,rounded corners=.25em](wh1.south)--+(0,-1em)-|(wh2.south)node [anchor=center,pos=0.25,fill=white] {\small\textit{im VF statt im MF}};
	\end{tikzpicture}
\z 

\vspace{1.25em}
\endgroup

Wenn das Vollverb wie \textit{getroffen} nicht finit ist, dann steht es auch im Verberst- und im Verbzweitsatz in seiner ursprünglichen Endstellung, d.\,h., es steht nach den von ihm abhängigen Einheiten. 
 
Betrachten wir das Mittelfeld, so zeigt sich die Eigenschaft des Deutschen, eine\is{Mittelfeld!Reihenfolge} Sprache mit freier Wortstellung zu sein. Die Konstituenten im Mittelfeld können wir frei umstellen. Diese im Abschnitt~\ref{subsection:Mittelfeld} dargestellte Möglichkeit bezeichnet man als \textit{Scrambling} (\gq{Durcheinanderbringen}). Die Umstellungsmöglichkeiten von \textit{gestern} in (\ref{ea:Vletztdeutsch}) illustrieren wir in den Sätzen in (\ref{ea:MFScrambling}).

\ea \label{ea:MFScrambling}
\ea \label{ea:MFscrambling1} dass die große Frau den kleinen Mann \textit{gestern} getroffen hat
\ex \label{ex:MFscrambling2} dass die große Frau \textit{gestern} den kleinen Mann getroffen hat
\ex \label{ex:MFscrambling3} dass \textit{gestern} die große Frau den kleinen Mann getroffen hat
\z 
\z 

Im Englischen ist nur die Voranstellung der Angabe \textit{yesterday} vor das Verb bzw. vor die VP  (\ref{ea:keinMFscramblingEnglisch1}) oder die Nachstellung (\ref{ex:keinMFScramblingEnglAdv}) möglich. Andere Stellungen wie (\ref{ex:keinMFscramblingEnglisch2}) sind nicht möglich.

\newpage
\ea \label{ea:keinMFScramblingEnglisch}
\ea [] {that the big woman \textit{yesterday} met the little man}\label{ea:keinMFscramblingEnglisch1}
\ex [] {that the big woman met the little men \textit{yesterday}}\label{ex:keinMFScramblingEnglAdv}
\ex [*] {that the little man met yesterday the big woman}\label{ex:keinMFscramblingEnglisch2}
\z 
\z 

Im Deutschen können wir wie in (\ref{ea:MFscramblingDeutsch}) die Ergänzungen \textit{die kleine Frau} und \textit{den kleinen Mann} umstellen, ohne die Bedeutung zu verändern. Im Englischen verändert sich bei solch einer Umstellung die Bedeutung. In (\ref{ex:keinMFscramblingEnglisch3}) bezeichnet \textit{the little man} nicht denjenigen, den \textit{the big woman} getroffen hat, sondern umgekehrt: der kleine Mann hat die große Frau getroffen. 

\ea
\ea [] {dass den großen Mann die kleine Frau getroffen hat}\label{ea:MFscramblingDeutsch}
\ex [] {that the little man met the big woman}\label{ex:keinMFscramblingEnglisch3}
\z 
\z 

Die relativ freie Wortstellung bezieht sich im heutigen Deutsch auf die Abfolge der verbalen Ergänzungen und Angaben. Die Stellung der verbalen Teile ist durch die Klammerstruktur auf der einen Seite und die Verbletztstellung auf der anderen Seite festgelegt. Somit ist es kein Widerspruch, vom Deutschen als Verbzweitsprache mit Verbletztstellung und freier Wortfolge zu sprechen.

\section{Zusammenfassung Feldermodell}\label{section:Zfsg_Feldermodell}

Bevor wir uns im folgenden Kapitel~\ref{chap:Konstituenz} mit den unmittelbaren Satzbausteinen und ihrer Ermittlung beschäftigen, fassen wir die wesentlichen Punkte zum Thema Feldermodell in Stichpunkten zusammen:

\begin{itemize}

\item Verbstellung: Verschiedene Stellungen des finiten Verbs erlauben die formale Differenzierung von drei Satztypen: Verbzweit-, Verberst- und Verbletztsätze.  

\item Satzarten: Den formalen Satztypen lassen sich Satzarten zuordnen, die sich durch die mit ihnen typischerweise vollzogenen sprachlichen Handlungen (Illokution) unterscheiden: Verbzweitsätze dienen als Aussagesätze dem Behaupten und als W-Fragesätze dem Erfragen von Entitäten. Mit Verberstsätzen erfragen wir das Zutreffen einer Aussage oder fordern zu deren Ausführung auf. Mit Verbletztsätzen handeln wir typischerweise nicht direkt im Sinne eines Aussagens, Erfragens oder Aufforderns.   

\item Satzklammer 1: Verbalkomplexe bilden in Verbzweit- und Verberstsätzen durch die getrennte Stellung des finiten Verbs und der infiniten Verbteile eine Klammer. Das finite Verb bildet die linke Satzklammer und die infiniten Bestandteile die rechte Satzklammer. Wenn das Verb einteilig und deshalb die rechte Satzklammer nicht besetzt ist, kann es \zB mit einem Modalverb zweiteilig gemacht werden, um die rechte Satzklammer sichtbar zu machen.

\item Satzklammer 2: In Verbletztsätzen kann ebenso von einer Klammer zwischen der einleitenden Subjunktion (linke Satzklammer oder Subjunktionsposition) und dem gesamten Verbalkomplex mit Endstellung des Finitums (rechte Satzklammer oder Verbalkomplex) gesprochen werden. 

\item Felder: Durch die Satzklammer entstehen Felder: 

\begin{itemize}

\item Mittelfeld: In allen Satztypen steht zwischen linker und rechter Satzklammer das Mittelfeld: Es kann mit beliebig vielen Konstituenten besetzt werden, deren Reihenfolge durch verschiedene Prinzipien gesteuert wird. Typischerweise stehen neue Informationen im Mittelfeld.

\item Vorfeld: In Verbzweitsätzen steht vor der linken Satzklammer das Vorfeld. Es muss mit einer Konstituente besetzt werden. Typischerweise steht das Thema der Äußerung oder der Rahmen (Ort/Zeit) der Aussage im Vorfeld. Vor dem Vorfeld gibt es für Linksversetzungen das Vor-Vorfeld. 

\item Nachfeld: In allen Satztypen steht das Nachfeld nach der rechten Satzklammer. Typischerweise wird es von Nebensätzen besetzt und dient der Entlastung des Mittelfelds. Die Nachfeldbesetzung durch Infinitive bzw. Infinitivgruppen ist nur dann möglich, wenn der Infinitiv bzw. die Infinitivgruppe satzwertig ist (kohärente vs. inkohärente Konstruktion).

\item Anschlussfeld: In allen Satztypen gibt es ganz links ein Anschlussfeld. Im Anschlussfeld können Konjunktionen stehen, die gleichrangige Sätze miteinander verbinden.

\end{itemize}

\item Feldermodell: Entweder beschreibt man die Stellungsregularitäten durch ein einheitliches Feldermodell (uniform) oder man benutzt für jeden Satztyp ein spezifisches Modell (differenziert). Das uniforme Modell zielt auf Gemeinsamkeiten der Satztypen und die Ableitung von Verberst- und Verbzweit- aus Verbletztsätzen ab. Das differenzierte Modell zielt auf die Unterschiede der einzelnen Satztypen ab.

\item Sprachtyp Deutsch: Bezüglich grundlegender Rektionsverhältnisse ist das Deutsche eine Verbletztsprache. Die Zweitstellung des Finitums (Verbzweitsatz) ist der häufigste Stellungstyp. Rein formal lassen sich aus der Endstellung des Finitums (Verbletztsatz) am einfachsten Verberst- und Verbzweitsätze ableiten. Die Stellungsmöglichkeiten im Mittelfeld sind relativ frei.

\end{itemize}

Der folgende letzte Abschnitt dieses Kapitels bietet Ihnen die Gelegenheit, durch Übungen zu lernen.


\section{Aufgaben zum Feldermodell}\label{section:Aufgaben_Feldermodell}

In den folgenden vier Aufgaben können Sie Ihr Wissen zum Thema Feldermodell anwenden. Wir empfehlen, dass Sie die Aufgaben zuerst selbstständig lösen. Im Anschluss können Sie Ihre Lösungen mit unseren Lösungen und Kommentaren vergleichen.

\subsection{Beschreibung nach dem Feldermodell}\label{subsection:FM1}

Wir haben Sätze nach der Stellung des finiten Verbs in Verbzweit- (vgl. Abschnitt~\ref{section:Verbzweitsätze}), Verberst- (vgl. Abschnitt~\ref{section:Verberstsatz}) und Verbletzt-Sätze (vgl. Abschnitt~\ref{section:Verbletztsatz}) unterschieden. In (\ref{ea:FM1_Sätze}) findet sich für jeden Typ ein Vertreter. 

\ea \label{ea:FM1_Sätze}
\ea \label{ea:FM1_V1} Kann, was gut ist, jemals vergessen werden?\footnote{Uwe Johnson: \textit{Jahrestage}. Bd. 1. Frankfurt a. M.: Suhrkamp 1970, S. 111.}
\ex \label{ex:FM1_VLetzt} Ob Anita bemerkte, daß sie die einzige war, die einen Bubikopf hatte?\footnote{Martin Walser: \textit{Ein springender Brunnen}. Frankfurt a. M.: Suhrkamp 1998, S. 146.}
\ex \label{ex:FM1_V2} Dass die Else, die dumme kleine Else, diesen gutbesoldeten älteren Eisenbahner zum Mann bekommen hatte, konnte ihre Cousine Auguste Marnet auch jetzt noch nicht verdauen.\footnote{Anna Seghers: \textit{Das siebte Kreuz}. Berlin: Aufbau-Taschenbuch 2002 [1942], S. 403.}
\z 
\z  

Nennen Sie für die Sätze in (\ref{ea:FM1_Sätze}) den Satztyp und zeigen Sie für diesen die Besetzung der Felder im uniformen oder im differenzierten Feldermodell (vgl. Abschnitt~\ref{section:differenziert_uniform_Feldermodell}).


Kommen wir nun zu den Lösungen: Der Satz in (\ref{ea:FM1_V1}) ist ein Verberstsatz. Das finite Verb \textit{kann} steht an erster Stelle. Das Finitum bildet genauso wie in Verbzweitsätzen die linke Satzklammer. Die rechte Satzklammer wird von den infiniten Verbteilen \textit{vergessen werden} gebildet. Zwischen linker und rechter Satzklammer spannt sich das Mittelfeld auf und ist hier durch \textit{was gut ist, jemals} besetzt. Das Nachfeld ist unbesetzt. Im uniformen Modell gibt es vor der linken Satzklammer ein Vorfeld, das aber obligatorisch unbesetzt bleibt. Im differenzierten Modell gibt es in Verberstsätzen kein Vorfeld. Das Anschlussfeld, das Konjunktionen besetzen können, ist unbesetzt. Wir zeigen die Besetzung im differenzierten Modell in der Abbildung~\ref{fig:FM1}. Die Besetzung der Felder des uniformen Modells zeigen wir zum Schluss für alle drei Sätze in der Abbildung~\ref{fig:FM5}.

\begin{figure}[!htbp]
  \centering
  %\resizebox{\textwidth}{!}{  
  \oneline{
    \begin{tabular}{cp{0.1em}cp{0.1em}cp{0.1em}cp{0.1em}c}
      AF && LSK && MF && RSK && NF \\
      \cmidrule{1-1}\cmidrule{3-3}\cmidrule{5-5}\cmidrule{7-7}\cmidrule{9-9}
       -- && Kann, && was gut ist, jemals && vergessen werden? && -- \\
    \end{tabular}
  }
  \caption{Lösung für Satz (\ref{ea:FM1_V1}) im differenzierten Modell}
  \label{fig:FM1}
\end{figure}   


Der Satz in (\ref{ex:FM1_VLetzt}) ist ein Verbletztsatz. Verbletztsätze bilden eine Klammer zwischen Subjunktion, hier \textit{ob}, und dem Verbalkomplex, an dessen Ende das finite Verb steht. In unserem Beispiel haben wir nur ein finites Verb, nämlich \textit{bemerkte}. Zwischen Subjunktion und Finitum bzw. Verbalkomplex mit Finitum spannt sich das Mittelfeld auf. In unserem Beispiel ist es durch \textit{Anita} besetzt. Die Nebensätze \textit{daß sie die einzige war, die einen Bubikopf hatte} besetzen das Nachfeld. Die den Nebensatz eröffnende Subjunktion bildet im uniformen Modell die linke Satzklammer und das Finitum (mit Verbalkomplex, falls vorhanden) die rechte. Im differenzierten Modell benennen wir die Stelle für die Subjunktion als Subjunktion und die Stelle für den Verbalkomplex als Verbalkomplex. Das Anschlussfeld ist in unserem Beispiel unbesetzt. Wir zeigen die Besetzung im differenzierten Modell in der Abbildung~\ref{fig:FM2}.

\begin{figure}[!htbp]
  \centering
  %\resizebox{\textwidth}{!}{
  \oneline{
    \begin{tabular}{cp{0.1em}cp{0.1em}cp{0.1em}cp{0.1em}c}
      AF && Subjunktion && MF && VK && NF \\
      \cmidrule{1-1}\cmidrule{3-3}\cmidrule{5-5}\cmidrule{7-7}\cmidrule{9-9}
      -- && Ob && Anita && bemerkte, && daß [\ldots{}]? \\

    \end{tabular}
  }
  \caption{Lösung für Satz (\ref{ex:FM1_VLetzt}) im differenzierten Modell}
  \label{fig:FM2}
\end{figure}
 
 \largerpage
Der Satz in (\ref{ex:FM1_V2}) ist ein Verbzweitsatz. Das Finitum \textit{konnte} steht an der zweiten Stelle und bildet die linke Satzklammer als Finitheitsposition. Verbzweitsätze zeichnen sich dadurch aus, dass das Vorfeld vor dem Finitum von einer Konstituente bzw. von einem Satzbaustein besetzt sein muss. In unserem Fall besetzt der Nebensatz \textit{dass die Else, die dumme kleine Else, diesen gutbesoldeten älteren Eisenbahner zum Mann bekommen hatte} das Vorfeld. Die rechte Satzklammer bildet der Infinitiv \textit{verdauen}. Zwischen der rechten und linken Klammer spannt sich das Mittelfeld auf. Es ist besetzt durch \textit{ihre Cousine Auguste Marnet auch jetzt noch nicht}. Das Nachfeld ist unbesetzt. Ebenfalls unbesetzt ist das Anschlussfeld. Zwischen dem uniformen und dem differenzierten Modell gibt es keine wesentlichen Unterschiede. Wir zeigen die Besetzung der Felder in der Abbildung~\ref{fig:FM4}.

\begin{figure}[!htbp]
  \centering
   %\resizebox{\textwidth}{!}{
   \oneline{
  \begin{tabular}{cp{0.1em}cp{0.1em}cp{0.1em}cp{0.1em}cp{0.1em}c}
    AF && VF && LSK && MF && RSK && NF \\
    \cmidrule{1-1}\cmidrule{3-3}\cmidrule{5-5}\cmidrule{7-7}\cmidrule{9-9}\cmidrule{11-11}
    -- && Dass die Else [\ldots{}], && konnte && ihre Cousine [\ldots{}] && verdauen. && -- \\
  \end{tabular}
  }
  \caption{Lösung für Satz (\ref{ex:FM1_V2}) im Feldermodell}
  \label{fig:FM4}
\end{figure}

Zum Schluss zeigen wir in der folgenden Abbildung~\ref{fig:FM5} die Besetzung der Felder im uniformen Modell.  

\begin{figure}[!htbp]
  \centering
   %\resizebox{\textwidth}{!}{
   \oneline{
  \begin{tabular}{cp{0.1em}cp{0.1em}cp{0.1em}cp{0.1em}cp{0.1em}c}
    AF && VF && LSK && MF && RSK && NF \\
    \cmidrule{1-1}\cmidrule{3-3}\cmidrule{5-5}\cmidrule{7-7}\cmidrule{9-9}\cmidrule{11-11}
    
    -- && -- && Kann, && was gut ist, jemals && vergessen werden && -- \\

    -- && -- && Ob && Anita && bemerkte, && daß [\ldots{}]\\
    
    -- && Dass die Else [\ldots{}], && konnte && ihre Cousine [\ldots{}] && verdauen. && -- \\
  \end{tabular}
  }
  \caption{Lösungen im uniformen Feldermodell}
  \label{fig:FM5}
\end{figure}


\subsection{Mittel- vs. Nachfeld}\label{subsection:FM2}
\largerpage
Bei Verberst- und Verbzweitsätzen kann die rechte Satzklammer (RSK) unbesetzt bleiben. Dann ist es, wie bei den beiden Sätzen in (\ref{ea:FM2_RSK_leer}), nicht ganz einfach, das Mittel- (MF) vom Nachfeld (NF) abzugrenzen.

\ea \label{ea:FM2_RSK_leer}
\ea \label{ea:FM2_RSK_leer1} Man glaubt, er sei ein übernatürliches lebendiges Wesen, das durch Geschenke besänftigt werden kann bzw. Wünsche erfüllt.\footnote{Walter Moers: \textit{Die 13 1/2 Leben des Käpt'n Blaubär}. Frankfurt a. M.: Eichborn 1999, S. 342.}
\ex \label{ex:FM2_RSK_leer2} Die glaubt den Hasen schon in der Pfanne.\footnote{Willi Bredel: \textit{Die Väter}. In: ders., Gesammelte Werke in Einzelausgaben. Bd. 7. Berlin: Aufbau-Verl. 1973 [1946], S. 236.}
\z 
\z 

Begründen Sie die jeweilige Mittel- und Nachfeldbesetzung.
Zeigen Sie für beide Sätze die Besetzung aller Felder in tabellarischer Form auf.


Kommen wir nun zu den Lösungen: Wenn es, wie in (\ref{ea:FM2_RSK_leer1}) und (\ref{ex:FM2_RSK_leer2}), keine infiniten Verbformen gibt, dann ist nicht ganz klar, wo die unbesetzte RSK zu verorten ist. In diesem Fall bilden wir aus dem einteiligen Verb ein zweiteiliges. Also statt \textit{glaubt} bilden wir \zB das Perfekt \textit{hat geglaubt}. Nun ist die RSK durch den infiniten Verbteil \textit{geglaubt} besetzt und die Mittel- bzw. Nachfeldbesetzung wird wie in (\ref{ea:FM2_RSK_besetzt}) deutlich.

\ea \label{ea:FM2_RSK_besetzt}
\ea [] {Man hat geglaubt, er sei ein übernatürliches lebendiges Wesen, das durch Geschenke besänftigt werden kann bzw. Wünsche erfüllt.}\label{ea:FM2_1}
\ex  [?] {Man hat, er sei ein übernatürliches lebendiges Wesen, das durch Geschenke besänftigt werden kann bzw. Wünsche erfüllt, geglaubt.}\label{ex:FM2_2}
\z 
\z 

Die fragliche Akzeptabilität von (\ref{ex:FM2_2}) spricht im Vergleich zu der vollkommen akzeptablen Stellung in (\ref{ea:FM2_1}) dafür, dass das NF durch \textit{er sei ein übernatürliches lebendiges Wesen, das durch Geschenke besänftigt werden kann bzw. Wünsche erfüllt} besetzt ist. Das MF ist unbesetzt.

Für den Satz in (\ref{ex:FM2_RSK_leer2}) fällt der Test hingegen anders aus, was beim Vergleich von (\ref{ea:FM2_3}) mit (\ref{ex:FM2_4}) ersichtlich wird. 

\ea \label{ea:FM2_RSK_besetzt_2}
\ea [*] {Die hat geglaubt den Hasen schon in der Pfanne.}\label{ea:FM2_3}
\ex [] {Die hat den Hasen schon in der Pfanne geglaubt.}\label{ex:FM2_4}
\z 
\z 

Da (\ref{ea:FM2_3}) im Gegensatz zu (\ref{ex:FM2_4}) ungrammatisch ist, besetzt \textit{den Hasen in der Pfanne} das MF. Das NF hingegen ist unbesetzt. Anders sähe die Besetzung aus, wenn der Satz lautete: \textit{Die glaubte, den Hasen schon in der Pfanne zu haben}. Dann würde \textit{den Hasen schon in der Pfanne zu haben} auch im Nachfeld stehen.

Die unterschiedlichen Feldbesetzungen stellen wir in der Abbildung~\ref{fig:FM3} dar. Der Übersichtlichkeit halber vernachlässigen wir das unbesetzte Anschlussfeld in der Darstellung.

\begin{figure}[!htbp]
  \centering
 % \resizebox{\textwidth}{!}{
 \oneline{
    \begin{tabular}{cp{0.1em}cp{0.1em}cp{0.1em}cp{0.1em}c}
      VF && LSK && MF && RSK && NF \\
      \cmidrule{1-1}\cmidrule{3-3}\cmidrule{5-5}\cmidrule{7-7}\cmidrule{9-9}
      Man && glaubt && -- && -- && er sei [\ldots{}].  \\
       Die && glaubt && den Hasen schon in der Pfanne. && -- && -- \\
    \end{tabular}
  }
  \caption{Lösung für die Sätze (\ref{ea:FM2_RSK_leer1}) und (\ref{ex:FM2_RSK_leer2}) im Feldermodell.}
  \label{fig:FM3}
\end{figure}

 
\subsection{Fehlerbeschreibungen}\label{subsection:FM3}
\largerpage
Die folgenden Äußerungen in (\ref{ea:FM_Fehler_DaF}) stammen von Lernenden von Deutsch als Fremdsprache.\footnote{Die Daten stammmen aus dem fehlerannotierten Lernerkorpus des Deutschen als Fremdsprache \textit{Falko}. Satz (\ref{ea_FM3_falsch}) unter \url{https://korpling.german.hu-berlin.de/falko-suche/?id=7cf68d6c-b68c-49cb-b177-3b2d8127cf28}, Satz (\ref{ex_FM4_falsch}) unter \url{https://korpling.german.hu-berlin.de/falko-suche/?id=50132c02-f52d-40fe-acde-f5fc78b732ea}, Satz (\ref{ex_FM5_falsch}) unter \url{https://korpling.german.hu-berlin.de/falko-suche/?id=841a5e32-a217-4611-a339-9702de34228c}.} Sie enthalten Stellungsfehler bzw. Abweichungen von der Standardschriftlichkeit, die sich mit Hilfe des Feldermodells beschreiben lassen. 

\ea \label{ea:FM_Fehler_DaF}
\ea [*] {Jetzt für Frauen ist es ganz anders in unsere gesellschaft.}\label{ea_FM3_falsch}
\ex [*] {Bevor ich anfange aber, möchte ich einige Hauptunterschiede zwischen dem EU und der Gründung der Vereinigten Staaten umreißen.}\label{ex_FM4_falsch}
\ex [*] {Obwohl auf den ersten Blick das Studium scheint nicht viel Wert zu sein, [\ldots{}]}\label{ex_FM5_falsch}
\z 
\z 

Beschreiben Sie nur diejenigen Stellungsfehler, die sich mit Hilfe des Feldermodells analysieren lassen, und korrigieren Sie diese. 


Kommen wir nun zu den Lösungen: In Satz (\ref{ea_FM3_falsch}) besteht der Stellungsfehler darin, dass zwei Konstituenten, zwei unmittelbare Satzbausteine, im Vorfeld stehen, nämlich \textit{jetzt} und für \textit{für Frauen}. In der Standardschriftlichkeit steht im deutschen Verbzweitsatz jedoch nur eine Konstituente, nur ein Satzbaustein, im Vorfeld, weil das finite Verb in der linken Satzklammer an zweiter Stelle steht. Diese Abweichung kann wie in (\ref{ea:FM1_korrekt}) oder in (\ref{ex:FM2_korrekt}) berichtigt werden.

\ea 
\ea \label{ea:FM1_korrekt} Jetzt ist es für Frauen ganz anders in unserer Gesellschaft.
\ex \label{ex:FM2_korrekt} Für Frauen ist es jetzt ganz anders in unserer Gesellschaft.
\z 
\z 

In Satz (\ref{ex_FM4_falsch}) liegt der Stellungsfehler im vorangestellten Verbletztsatz. In Verbletztsätzen steht das Finitum an letzter Stelle. Die Konjunktion \textit{aber} kann nicht im Nachfeld stehen. Sie muss entweder wie in (\ref{ea:FM3_korrekt}) im Anschlussfeld stehen oder wie in (\ref{ex:FM4_korrekt}) im Mittelfeld.

\ea 
\ea \label{ea:FM3_korrekt} Aber bevor ich anfange, möchte ich einige Hauptunterschiede zwischen der EU und der Gründung der Vereinigten Staaten umreißen.
\ex \label{ex:FM4_korrekt} Bevor ich aber anfange, möchte ich einige Hauptunterschiede zwischen der EU und der Gründung der Vereinigten Staaten umreißen.
\z 
\z 

Bei Satz (\ref{ex_FM5_falsch}) handelt es sich um einen Verbletztsatz. Der Verbalkomplex besteht aus dem Modalitätsverb \textit{scheint} (vgl. Abschnitt~\ref{subsection:Modalitätsverb_zu_Inf}) und dem von ihm regierten \textit{zu}-Infinitiv \textit{zu sein}. Modalitätsverben konstruieren in der Regel kohärent. In der rechten Satzklammer bzw. im Feld für den Verbalkomplex steht zuerst der \textit{zu}-Infinitiv und danach, wie in (\ref{ea:FM5_korrekt}), das ihn regierende Modalitätsverb (vgl. Abschnitt~\ref{subsection:KohärenteKonstruktion}). Eine inkohärente Konstruktion mit dem \textit{zu}-Infinitiv im Nachfeld (\ref{ea:FM6_inkorrekt}) ist normalerweise nicht möglich.  

\ea 
\ea [] {Obwohl auf den ersten Blick das Studium nicht viel wert zu sein scheint, [\ldots{}].}\label{ea:FM5_korrekt}
\ex [*] {Obwohl auf den ersten Blick das Studium scheint nicht viel wert zu sein, [\ldots{}].}\label{ea:FM6_inkorrekt}
\z 
\z 


\subsection{Verbalkomplex}\label{subsection:FM4}

In den beiden folgenden Verbletztsätzen sind die Stellungen der Verbteile innerhalb der rechten Satzklammer bzw. innerhalb des Felds für den Verbalkomplex unterschiedlich.

\ea \label{ea:FM_Verbalkomplex}
\ea \label{ea:FM_Verbalkomplex_Rektionsrichtung} Kim war die erste Person, abgesehen von meiner Mutter, die mich singen gehört hat.\footnote{Süddeutsche Zeitung, 05.07.2008.}
\ex \label{ex:FM_Verbalkomplex_OF_UF} Wer ihn je hat singen hören, wünscht sich, er sänge nicht.\footnote{Süddeutsche Zeitung, 18.02.2006.}
\z 
\z  

Beide Abfolgen, \textit{singen gehört hat} in (\ref{ea:FM_Verbalkomplex_Rektionsrichtung}) und \textit{hat singen hören} in (\ref{ex:FM_Verbalkomplex_OF_UF}), gelten standardsprachlich als korrekt. Wie lassen sich die beiden Stellungen im Rahmen des Feldermodells beschreiben? Denken Sie an die Rektionsrichtung und die Unterscheidung zwischen Ober- und Unterfeld innerhalb der rechten Satzklammer bzw. des Felds für den Verbalkomplex bei Verbletztsätzen (vgl. Abschnitt~\ref{subsection:RechteSatzklammer}).


\largerpage
Kommen wir nun zu den Lösungen: Die Abfolge in (\ref{ea:FM_Verbalkomplex_Rektionsrichtung}) folgt der Regel, dass in Verbletztsätzen die Abfolge der verbalen Teile umgekehrt zu den Rektionsverhältnissen verläuft. Das finite Verb an letzter Stelle \textit{hat} regiert das vor ihm stehende Partizip II \textit{gehört}. Dieses regiert wiederum den vor ihm stehenden Infinitiv \textit{singen}. Diese Abfolge -- entgegengesetzt zur Rektionsrichtung \textit{singen gehört hat} --gilt für das Unterfeld innerhalb der rechten Satzklammer bzw. des Felds für den Verbalkomplex in Verbletztsätzen.

Wenn es jedoch, wie in unserem Fall, neben dem finiten Verb noch zwei infinite Bestandteile gibt, so kann das Finitum, wie in (\ref{ex:FM_Verbalkomplex_OF_UF}), an der Spitze des Verbalkomplexes stehen. Diese Position haben wir als Oberfeld bezeichnet. Nun heißt es aber standardsprachlich nicht \textit{hat singen gehört}, sondern \textit{hat singen hören}. Die Form des im Unterfeld stehenden Partizip II hat sich verändert. Statt \textit{gehört} steht dort \textit{hören}. Der Infinitiv wird als Ersatz für das Partizip II gebraucht, weshalb man von einem Ersatzinfinitiv spricht.    

